\chapter{Cronología (1690--2026)}
\label{cap:cronologia}

Entradas marcadas con asterisco: documentadas en fuente oficial o judicial con fecha exacta verificada. Cada entrada de 1908 a 1916 se desarrolla, con su número de Boletín y su hoja, en el capítulo correspondiente o en el apéndice \ref{cap:dominio}; para las posteriores, el lugar donde se acredita es el capítulo que la trata. Las entradas sin asterisco anteriores a 1990 provienen igualmente del Boletín Oficial y se desarrollan en el capítulo correspondiente; las posteriores sin asterisco provienen de cobertura periodística o de partes oficiales, y se consignan como tales.

\begin{longtable}{@{}p{2.7cm}p{10.3cm}@{}}
\toprule
\textbf{Fecha} & \textbf{Hecho y fuente} \\
\midrule
\endhead
1690 & Francisco Vaquero, poseedor de las tierras, desarrolla agricultura, ganadería y molienda. Origen del topónimo. \\
1735* & Registro del molino hidráulico de piedra conservado en el museo popular de Vaqueros. \\
18 nov.\ 1866* & \textbf{Ordenanza municipal de Campo Santo} que reparte el agua del río Mojotoro; la Legislatura la sanciona como ley en enero de 1867. Es el antecedente más antiguo del reparto de esta cuenca y funda los derechos que en 1908 se defenderán en juicio. Se conoce por el fallo de 1910 y no fue contrastada contra la norma. \\
28 nov.\ 1884* & \textbf{Ley 319}: autoriza a la Municipalidad de La Caldera a percibir derechos de inhumación ---un peso por adulto, cincuenta centavos por párvulo, gratis para los pobres de solemnidad--- y la obliga a llevar un libro de defunciones. Es la mención legislativa más antigua del municipio que este relevamiento registra. \\
1885 & \textbf{La Municipalidad de La Caldera concede} a Eustaquio Murúa el uso de treinta litros por segundo del río Wierna para regar la finca «Sausal o Curuzú». \\
1908 & Aparece el Boletín Oficial de Salta. Pedimento de cateo de cobre en el ``partido del Potrero de Castillo'', en terrenos de «Garzon y Pinto, vecinos de Buenos Aires». Juicio Ruiz c/ Arancibia: se pierde el título de Bautista Wierna (la sentencia no consigna el departamento). En noviembre, jurados de reclamos y clasificador de patentes de Caldera para 1909; en diciembre, Escolástico Aparicio\index{Aparicio, Escolástico M.}, comisario de policía del departamento para 1909. \\
25 sep.\ 1908* & \textbf{Ley que dota de agua al pueblo de General Güemes}. Tres particulares la demandan por inconstitucional alegando derechos adquiridos sobre el caudal del \textbf{río Mojotoro} en virtud de una \textbf{ordenanza municipal de Campo Santo del 18 de noviembre de 1866}, sancionada como ley en enero de 1867. \textbf{El reparto del agua de esta cuenca está en disputa desde el primer año del Boletín.} \\
1909 & Cuadro de rentas municipales de 1908: el municipio \textbf{vende su recaudación en licitación}. El P.\ E.\ nombra por decreto a los \textbf{cinco miembros de la Comisión Municipal}, porque el sorteo de renovación no se había hecho; tres renuncian en el año y los reemplazan Figueroa\index{Figueroa, Abelardo}, Vidal y Regis\index{Regis, Augusto}. Decreto de comisarios: \textbf{trece partidos} del departamento; dos meses después, los \textbf{veedores del Código Rural} dividen el mismo territorio en \textbf{quince secciones}, entre ellas Wierna. Jueces de paz: Carlos Matienzo y Augusto Regis\index{Regis, Augusto}, que renuncia. \textbf{Licitación de un edificio para la Policía de La Caldera.} \\
1909 & \textbf{Deslinde general de las fincas San Alejo y Santa Rufina}, con incidentes que siguen en 1910. Ese mismo año se remata una \textbf{décima parte indivisa} de la finca Angosto de Arrieta y, dos veces, la sucesión Catacata en Potrero de Castillo. \\
1910 & Los vecinos piden defender el pueblo ``de las invasiones del Río''. La Provincia asigna \$2.000 y nombra comisión. \textbf{Escolástico Aparicio}\index{Aparicio, Escolástico M.} acumula en el año cuatro cargos ---miembro de la Comisión Municipal, comisario de policía del departamento, comisionado de empadronamiento y otra vez comisario y miembro---, y el decreto de diciembre repite el argumento de 1909 para designar las comisiones municipales sin sorteo. En junio, \textbf{doce partidos}: ya aparecen la estación Mojotoro y Angostura de Arrieta y Campo Alegre. En agosto se aprueba el \textbf{contrato con Moisés Vera para construir el edificio de la comisaría}. Caldera elige \textbf{un senador}. \\
1910 & Mensura de las fincas \textbf{Los Porongos} y \textbf{Loma Bola}. \\
12 feb.\ 1910* & El Superior Tribunal resuelve que la demanda contra la ley de aguas de General Güemes no es contencioso-administrativa sino ordinaria por inconstitucionalidad: «\textit{las partes pueden clasificar su acción, por cualquier causa, erróneamente, sin que esto importe alterar ó cambiar la naturaleza de la causa}». No resuelve el fondo. \\
1 ago.\ 1910* & La Provincia aprueba el contrato con Moisés Vera para la \textbf{construcción del edificio de la Comisaría de La Caldera} (B.O.\ Nº 178, 10/08/1910). No consta monto, plazo ni terreno. Treinta y tres años después el Estado deducirá posesión treintañal sobre el predio donde funciona la comisaría. \\
1911 & Sobrecosto de las defensas de 1910: \$958,65 más, sin autorización legislativa previa. El decreto de comisarios confirma el partido \textbf{Campo Alegre y Angosto de Arrieta}, que ya existía en 1910. Se jubila la directora de la \textbf{escuela elemental de la parroquia de la Caldera}. En octubre se crea una oficina del Registro Civil en la \textbf{estación Mojotoro}, ad honorem, por la distancia y la falta de transporte; en diciembre, un agente para la comisaría de \textbf{Vaqueros} por la obra del \textbf{puente} sobre ese río y por la \textbf{explotación de canteras} junto al puente del Mojotoro. En noviembre, \textbf{Escolástico Aparicio}\index{Aparicio, Escolástico M.} es comisionado del empadronamiento de capitales en giro y de la clasificación de patentes para 1912 (Nº 295, h.\ 2): cuarto año seguido con cargo. \textbf{Remate de San Alejo y Santa Rufina}: 17.874 ha, base \$180.000. \\
jul.\ 1911 & Jubilación de \textbf{Jacoba Luquez de Gómez}, directora de la «Escuela elemental de la parroquia de la Caldera», con \$78,74 mensuales (B.O.\ Nº 268, h.\ 3). \textbf{Primer nombre de la escuela del pueblo en este relevamiento} (capítulo \ref{cap:presencia}). \\
oct.\ 1911 & Se crea, \textit{ad honorem} y sólo para el enrolamiento, una oficina del \textbf{Registro Civil en la estación Mojotoro}, a cargo de \textbf{Benjamín López}, subcomisario del partido de la estación (B.O.\ Nº 291, h.\ 3; capítulo \ref{cap:presencia}). \\
1912 & \textbf{Segundo remate de San Alejo y Santa Rufina}, con la base rebajada de \$180.000 a \$135.000; el aviso declara una \textbf{Escuela Nacional} en la finca. Deslindes de \textbf{Los Porongos o San Luis} y de \textbf{Santa Mónica y Quitilipe}, las dos sobre el río La Caldera. Una sentencia ubica la finca Mojotoro «en jurisdicción del Departamento de La Caldera (no Campo Santo)». El padrón del departamento llega a \textbf{524} y se vota en la Escuela del Pueblo y el Juzgado de Paz. Se nombra \textbf{subcomisario del pueblo} con \$100 mensuales. Escolástico Aparicio\index{Aparicio, Escolástico M.} vuelve a ser comisario del departamento; los jueces de paz se nombran \textbf{de ternas de la Comisión Municipal}, y ésta se integra con Daniel Linares\index{Linares, Daniel} por renuncia de Augusto Rejis\index{Regis, Augusto}. \\
dic.\ 1912 & \textbf{Silverio Cáseres}, subcomisario interino del pueblo de La Caldera, con cien pesos mensuales; Benjamín López pasa a comisario del partido de Mojotoro (B.O.\ Nº 379, h.\ 3; capítulo \ref{cap:presencia}). \\
1913 & \textbf{Escolástico Aparicio}\index{Aparicio, Escolástico M.} renuncia a la comisaría del departamento en febrero y en abril es \textbf{juez de paz propietario}, propuesto en terna por la comisión municipal, que ese año son Daniel Linares\index{Linares, Daniel}, Camilo Cáceres y Salvador Rosa\index{Rosa, Salvador}. Trece partidos con comisario auxiliar. \textbf{Remate de la finca Lesser} ---«el punto más pintoresco de la provincia», a una hora en carruaje de la capital--- y de una \textbf{casa-quinta de siete habitaciones en el pueblo}. \\
1914 & El municipio de La Caldera dicta una \textbf{ordenanza que reglamenta la distribución de la totalidad del caudal del río Santa Rufina} y reconoce los «derechos adquiridos por uso tradicional». Texto no encontrado: no está en ninguna de las setenta ediciones de 1914 ni en las cuarenta y nueve de 1915, leídas todas una por una. \textbf{San Alejo y Santa Rufina sale a remate tres veces} ---21 de febrero, 12 de mayo y 5 de octubre--- con la base cayendo de \$5,50 a \$3,20 la hectárea. \textbf{Deslinde de la finca Campo Alegre}, pedido por Francisco Sosaya; agrimensor Ettore Chiostri (Nº 512). \textbf{Remate de la finca Angostura}, base \$20.000 (Nº 521). Escolástico Aparicio\index{Aparicio, Escolástico M.} vuelve a la comisión municipal en la vacante de Salvador Rosa\index{Rosa, Salvador}: séptimo año seguido con cargo. Comisario del departamento, Julio del Solar; subcomisarios de la estación Mojotoro, José Nogales y después Benjamín López; comisario auxiliar de Los Porongos, Carlos Bongarnier. El juez de paz del departamento firma el sucesorio de Salustiano Caro. \\
1915 & \textbf{Año leído edición por edición} (49 ediciones, falta la 581 y faltan enero y la primera semana de febrero). \textbf{26 de enero: decreto que nombra la comisión municipal} ---Eustaquio Murúa\index{Murúa, Eustaquio} y otra vez Adolfo Vidal, por haber terminado el período de Vidal y Ángel Fernández---. Caldera es convocada a elegir \textbf{senador y diputado} en el mismo artículo. \textbf{El padrón se reparte en dos mesas}: 1 a 298 en la Escuela Elemental Mixta del Pueblo y 1 a 216 ---o 316--- en la \textbf{Escuela Infantil Provincial de Vaqueros}. En mayo, un decreto de supresión de puestos «por razones de renta» \textbf{suprime el agente de la comisaría de campaña de Caldera}. El 2 de septiembre el \textbf{juez de paz E.\ M.\ Aparicio}\index{Aparicio, Escolástico M.} abre el sucesorio de Asunción Espeche González de Wayar. Tres partidos proclaman electores de gobernador por el departamento. En diciembre, \textbf{Aparicio}\index{Aparicio, Escolástico M.} es comisionado para clasificar las patentes de 1916 y los \textbf{jurados} de los reclamos son Daniel Linares\index{Linares, Daniel} y Silvano Murúa\index{Murúa, Silvano Ignacio}, con Adolfo Vidal y Belisario Nieva de suplentes. \textbf{Ningún acto de tierra, ningún balance municipal y ningún plano del departamento en todo el año.} \\
1916 & \textbf{Año leído edición por edición} (49 ediciones, 587 a 635, sin saltos y con las 196 hojas). \textbf{Tres actos de dominio}: deslinde de \textbf{Loma Bola} pedido por Alberto Escudero, perito Herman Pfister\index{Pfister, Hermann}; \textbf{remate judicial de Los Porongos hoy San Luis}, por ejecución de Enrique Werner contra el conde \textbf{Humberto Graf von Garnier Turawa}, base \$33.333,33 y \textbf{5.359 ha 77 a 83 ca}, la única superficie que un aviso del departamento declara en estos años; y \textbf{segundo deslinde de Santa Mónica y Quitilipe}, pedido por Manuel Goica, con una línea del ingeniero Pedro Cornejo del \textbf{9 de abril de 1889} como divisoria. \textbf{Escolástico Aparicio}\index{Aparicio, Escolástico M.} renuncia en febrero a la comisión municipal y en abril es \textbf{juez de paz propietario} por terna de esa misma comisión: noveno año seguido con cargo. \textbf{Eustaquio Murúa}\index{Murúa, Eustaquio} junta en cinco meses \textbf{comisario de policía}, \textbf{expendedor de guías} ---con fianza de Silvano I.\ Murúa\index{Murúa, Silvano Ignacio} por \$500--- y \textbf{receptor del impuesto al consumo}. La comisión municipal se renueva dos veces por renuncia: entra Andrés Marín, sale Adolfo Vidal, entra Salvador Rosa\index{Rosa, Salvador}. Tres mesas receptoras por insaculación. En noviembre, \textbf{pedido de cateo por plomos argentíferos en el Cerro Nevado}, terrenos de Garzón y Pintos, con dos pedimentos anteriores sobre la misma finca, uno de ellos de Casiano Hoyos \textbf{por sí y por Juan Larrán}\index{Larrán, Juan}. El 20 de febrero cambia el gobernador: Patrón Costas deja el mando y asume Cornejo. \textbf{Vaqueros no aparece una sola vez en todo el año.} \\
1917 & \textbf{Año leído edición por edición} (47 ediciones, 636 a 683; falta entera la 640 y 40 de las 329 hojas; sin OCR propio). \textbf{Tres deslindes}: el de la finca \textbf{Potrero Gallinato} ---con división de condominio, pedido por Camilo Cáceres, «que fue antes de Inocencio Aramayo», con \textbf{José Nogales} de lindero este, el mismo que en 1922 será el ejecutado cuya finca sale a remate---; el \textbf{deslinde parcial de un terreno en La Calderilla} de Juana Flores de Chaile, lindero de \textbf{Ildefonsa Burgos de Yugra} y de la familia Yugra; y el de unos terrenos del \textbf{partido de Calderilla} de Abraham Fernández y del menor Bernabé Fernández, con \textbf{Teodoro Reyes} de lindero norte. Perito de los tres, Herman Pfister\index{Pfister, Hermann}. El 4 de enero se \textbf{renueva entera la comisión municipal} (Daniel Linares\index{Linares, Daniel}, Andrés Marini, Camilo Cáceres) y el 23, por terna de ese mismo cuerpo, \textbf{Escolástico M. Aparicio}\index{Aparicio, Escolástico M.} es juez de paz propietario: décimo año seguido con cargo. \textbf{Eustaquio Murúa}\index{Murúa, Eustaquio} sigue de comisario de policía. En septiembre entra a la comisión municipal \textbf{Rafael Delgado}\index{Delgado, Rafael}, por renuncia de Marini. Linares es \textbf{candidato a diputado} por el departamento mientras integra su comisión municipal, y Caldera elige ese año un diputado y ningún senador. \textbf{Jubilación de Sara Blasco, directora de la Escuela Infantil de Vaqueros}, con \$42,24 mensuales. \textbf{Ningún plano en todo el año.} \\
1917 & Se constituye en Buenos Aires ``Levilly y Compañía'' para explotar la \textbf{fundición de Mojotoro} y una mina de plomo en Jujuy. \\
1918 & \textbf{Año leído edición por edición} (52 ediciones, 684 a 736; falta entera la 703 y 53 de las 587 hojas; sin OCR propio). \textbf{Ningún acto de dominio.} \\
4 mar.\ 1918* & \textbf{Decreto 1.557} en acuerdo de Ministros: convierte en \textbf{Municipalidades electivas}, desde el 1.\ de julio, a las Comisiones Municipales de catorce departamentos con renta superior a \$5.000. \textbf{Caldera está entre ellos.} \\
7 mar.\ 1918* & \textbf{Decreto 1.564}: nombra los comisarios auxiliares de los \textbf{trece partidos} del departamento a propuesta del comisario departamental. Vuelve Liborio Guerrero\index{Guerrero, Liborio} a Peñones y Chalchami; José Nogales toma Potrero de Gallinato y renuncia en junio. \\
20 abr.\ 1918* & \textbf{Decreto 1.653}: jubila a \textbf{Juana Luque}, Directora de la Escuela de la Caldera, con \$104,50 mensuales. Segundo nombre de la escuela del pueblo en este relevamiento, después de la directora jubilada en 1911 (capítulo \ref{cap:presencia}). \\
24 abr.\ 1918* & \textbf{Decreto 1.660}: deja sin efecto el 1.557 \textbf{sólo para Caldera}. El cálculo de recursos de \$6.471,44 incluía \$3.083,94 de saldo del año anterior: la renta real era \$3.387,50 y el departamento «\textit{no está en condiciones de ser administrado por una Municipalidad electiva}». Impreso «Abril 24 de 1913» [\textit{sic}]. \\
6 may.\ 1918* & \textbf{Decreto 1697}: jueces de paz del departamento, propietario \textbf{Casimiro Caseres} y suplente Gervasio Chalar, a propuesta en terna de la Comisión Municipal. \\
16--18 may.\ 1918* & La \textbf{Intervención Nacional} suspende a los trece comisarios auxiliares \textit{ad-honorem} del departamento (decreto 27) y manda un oficial de línea, el Tte.\ \textbf{Horacio Hornstein} (decreto 39). La numeración de los decretos reinicia en 1. \\
21 may.\ 1918* & \textbf{Decreto 50}: el Interventor deja sin efecto el 1.557 \textbf{entero} y suspende la inscripción de electores. Ninguno de los catorce departamentos eligió su municipalidad. \\
20 jun.\ 1918* & \textbf{Decreto 157}: subdivide la provincia en distritos mineros para el registro gráfico. Caldera queda en el \textbf{V, Distrito del Toro}. El decreto no acompaña el gráfico que manda llevar. \\
12 jul.\ 1918* & Pedimento de la mina de plomo «\textbf{20 de Febrero}», finca «Cerro Nevado» de Garzón y Pintos, por Miguel A.\ Sosa y Dermidio Arancibia. En agosto, los mismos denuncian el despueble de la mina «Sara». \\
23 sep.\ 1918* & \textbf{Decreto 113}: \textbf{jueces de agua} de Caldera y Calderilla, «\textit{vista la nota que antecede del Sr.\ Comisionado Municipal}». El decreto 190 del 7 de noviembre vuelve a nombrarlos con otros nombres y no cita al primero. \\
22 nov.\ 1918* & El \textbf{Colegio Electoral N.\ 3 --- Caldera} publica sus cuatro mesas y sus circuitos: \textbf{treinta y dos parajes}, siete de ellos nunca registrados antes por este libro. \\
1919 & \textbf{Año leído edición por edición} (49 ediciones, 738 a 786: los cuarenta y nueve viernes que van del 17 de enero al 19 de diciembre, de los cincuenta y dos del año; 440 hojas de 478 folios; las 38 hojas ausentes caen todas al final de veintiocho ediciones y ninguna en el medio). \textbf{Ningún acto de dominio, ningún plano, ningún resumen de Tesorería.} \\
15 ene.\ 1919* & \textbf{Decreto 33}: nombra la Comisión Municipal de Caldera ---\textbf{Domingo Royo, Francisco Sosaya, José Molina, Julio Mercado y Camilo Cáceres}--- y funda la decisión en que el art.\ 16 de la Ley Orgánica «\textit{ha sido derogado por el art.\ 175 de la Constitución, que exije para la existencia de Municipalidades autónomas un centro urbano de cinco mil habitantes, de que carecen todos los Departamentos de Campaña}». \textbf{El umbral ya no es la renta de 1918: es la población.} \\
4 feb.\ 1919* & \textbf{Decreto 117}: juez de paz propietario \textbf{Casimiro Cáceres} y suplente Carmelo Sivila, «\textit{vista las ternas presentadas por la Comisión Municipal}» [\textit{sic}]. Se publica diez días después. \\
7 mar.\ 1919* & \textbf{Decreto 213}, publicado ese día y \textbf{sin fecha propia impresa}: comisarios de \textbf{diez partidos} del departamento ---eran trece en 1909 y en 1918---. Aparecen como partido \textbf{El Sauce} y \textbf{el Monte}; Liborio Guerrero\index{Guerrero, Liborio} pasa de Peñones y Chalchami al Monte. \\
8 may.\ 1919* & \textbf{Decreto 337}: acepta la renuncia de Liborio Guerrero\index{Guerrero, Liborio} a los partidos de «\textit{Monte y Chalchalnio}» y nombra a Héctor Echanique. Cuarta grafía del mismo paraje. \\
11 jul.\ 1919* & \textbf{Decreto 414}: deja sin efecto el nombramiento de Pedro Vale en Lesser y nombra a Marcelo Lamas. El decreto no dice por qué. \\
19 sep.\ 1919* & \textbf{Ley de Presupuesto} del ejercicio de agosto a diciembre, publicada ese día, \$767.599,95 para cinco meses. Su ítem VII asigna a Caldera, entre cincuenta y cinco localidades, \textbf{un encargado de Registro Civil a \$44 cada uno «\textit{inclusive gastos de oficina y alquiler}»}. El ítem V, que paga a los comisarios de campaña, no cierra por sesenta pesos. \\
20 sep.\ 1919* & \textbf{Decreto 522}: crea la \textbf{Sub-Comisaría de Mojotoro «\textit{con sueldo}»} y nombra a Silverio Cáceres, con retroactividad al 1.\ de agosto. Es la primera plaza rentada de policía en un partido del departamento que este libro documenta, y \textbf{el acto no nombra al departamento}. \\
1920 & \textbf{Año leído edición por edición} (48 ediciones, 788 a 839; faltan la 793, la 822, la 837 y la 838; 526 hojas distintas sobre seiscientos folios). \textbf{Ningún acto de dominio, ningún plano, ningún balance municipal} ---y diez resúmenes de la Tesorería provincial, más que en ningún otro año con recuento. \\
4 feb.\ 1920* & \textbf{Decreto 677}: acepta la renuncia de \textbf{Silverio Cáceres} al puesto de Sub-Comisario de Mojotoro. No dice quién lo reemplaza. \\
6 feb.\ 1920* & \textbf{Decreto 685}: convoca a elecciones de legisladores provinciales para el 7 de marzo. \textbf{«Caldera un Diputado» figura dos veces en la misma enumeración}, que además nombra doce departamentos distintos en trece renglones y deja sin nombrar al menos ocho. No se concilia. Ubica las cuatro mesas del departamento en los mismos cuatro edificios de 1918. \\
26 feb.\ 1920* & \textbf{Decreto 724}: nombra a \textbf{José Agustín Martínez} Sub-Comisario de Mojotoro «\textit{para ocupar el puesto vacante}», sin mencionar la renuncia que el Boletín publicó tres semanas antes. \\
27 feb.\ 1920* & \textit{(fecha de la edición; las piezas no llevan fecha propia)} Tres partidos proclaman candidato a diputado por La Caldera: \textbf{José María Ocaranza} (U.C.R.\ Intransigente), \textbf{José Antonio Aráoz} (Unión Provincial) y \textbf{Dermidio Arancibia} (U.C.R.), que en 1918 fue uno de los dos peticionantes de las minas del departamento. Y la Junta Electoral reemplaza a \textbf{Augusto Regis}\index{Regis, Augusto} como presidente de la mesa 1, sin decir por qué. \\
22 mar.\ 1920* & \textbf{Decreto 772}: convoca a elecciones complementarias para el 4 de abril por \textbf{anulación de urnas}. La primera de la lista es «\textit{La Caldera --- C.\ Electoral N.º 3 Mesa N.º 4}», que es la de Mojotoro. El decreto no dice el lugar ni la causa. \\
16 jun.\ 1920* & \textbf{Decreto 901}: nombra a \textbf{Ciro Cáceres} Sub-Comisario de Mojotoro. \textbf{Tercer movimiento del mismo puesto en cinco meses}, y otra vez sobre una vacante declarada sin nombrar a quien la dejó. \\
25 oct.\ 1920* & \textbf{Decreto 1140}: nombra a \textbf{Casimiro Cáceres} ---juez de paz del departamento en 1918 y 1919--- encargado del Registro Civil de La Caldera, «\textit{con antigüedad al 5 de Agosto ppdo.}». \textbf{El interinato que reconoce nunca se publicó.} \\
26 nov.\ 1920* & \textit{(fecha de la edición; el decreto no lleva fecha impresa)} Decreto en acuerdo de Ministros: \textbf{Coroneles de la Guardia Nacional} por departamento. Por Caldera, \textbf{Mariano Coll}. El departamento no tiene teniente coronel. \\
1921 & \textbf{Año leído edición por edición} (43 ediciones efectivas de la 840 a la 888; seis marcadores de ausencia; 742 hojas; folios 1.681 a 2.532, que empalman sin hueco con el cierre de 1920). \textbf{Dos actos de dominio}, los primeros desde 1917; ningún plano; \textbf{ningún balance ni resumen de Tesorería}, después de los diez de 1920. \\
28 ene.\ 1921* & Edicto de \textbf{deslinde, mensura y amojonamiento de las fincas «Cerro de Buena Vista» y «Sausal o Curuzú»}, partido de Vaqueros, «\textit{de propiedad de don Francisco S.\ Urquiza}»\index{Urquiza, Francisco S.}; agrimensor Héctor Chiostri, títulos de fs.\ 33 y 40. Tres años después obtendrá cien litros por segundo del Wierna para la primera; la segunda lleva treinta desde 1885. \\
11 feb.\ 1921* & Decreto en acuerdo de Ministros: convoca a elecciones de legisladores para el 6 de marzo. Caldera elige \textbf{un diputado y un senador}. La U.C.R.\ proclama por el departamento al senador \textbf{Emeterio Royo} y al diputado \textbf{Lucio Ortiz}. \\
18 feb.\ 1921* & La Junta Electoral publica los \textbf{circuitos del Colegio Electoral Nº 3}: \textbf{treinta y cuatro parajes} en cuatro mesas, dos más que en 1918 ---aparecen \textbf{El Bordo} y la \textbf{Estación F.\ C.\ C.\ N.}---. Los cuatro edificios son los mismos desde 1918. \\
7 jul.\ 1921* & \textbf{Decreto 1669}: jueces de paz del departamento «\textit{vistas las ternas elevadas por la Comisión Municipal}», propietario \textbf{Laureano Peralta} y suplente \textbf{Escolástico M.\ Aparicio}\index{Aparicio, Escolástico M.}, que vuelve al departamento tres años después de su traslado a Campo Santo. El decreto no nombra a los salientes. \\
19 ago.\ 1921* & \textbf{Edicto de deslinde de la finca «El Angosto»}, de Antonia P.\ de Peralta. Sus linderos ---\textbf{Julio Mercado, Susana Aguilera} y un \textbf{Augusto Reyes}--- son los mismos tres que el remate del Angosto de Arrieta declaraba en 1909, doce años antes (apéndice \ref{cap:dominio}). \\
12 sep.\ 1921* & \textbf{Decreto 1816}: acepta la renuncia de Escolástico M.\ Aparicio\index{Aparicio, Escolástico M.} al juzgado de paz suplente. \textbf{Dos meses y cinco días}, y no se publica reemplazante. \\
20 sep.\ 1921* & \textbf{Decreto 1828}: la Junta Electoral \textbf{declara nula la elección del 7 de agosto en La Caldera} y en otros once distritos, «\textit{por no haberse verificado elección legal en dos tercios de las mesas}», y se convoca de nuevo para el segundo domingo de octubre. \textbf{La convocatoria se publica el 28 de octubre, diecinueve días después de la fecha señalada.} \\
7 y 11 nov.\ 1921* & \textbf{Decretos 2021 y 2062}: \textbf{Julio Royo} pierde en cuatro días el cargo de Comisario de Policía y el de Receptor de Rentas y Expendedor de Guías del departamento, y en los dos lo reemplaza un \textbf{Tedin}, con fianza de doña Encarnación Tedin de López por \$5.000. Ninguno de los dos decretos da la causa. \\
25 nov.\ 1921* & \textbf{Presupuesto de escuelas de la campaña}: el departamento tiene \textbf{dos escuelas} y cinco cargos ---la Elemental de 2.\textsuperscript{a} de la Parroquia, \$345 mensuales, y la Infantil de 2.\textsuperscript{a} de Vaqueros, \$210---. Primer dato de dotación escolar del departamento. \\
1922 & \textbf{Primer semestre leído edición por edición} (27 ediciones, 890 a 916, del 13 de enero al 21 de julio; 446 hojas; falta entera la 889, no hubo edición el viernes 14 de abril y la cadena de folios indica unas trece hojas ausentes dentro de las ediciones entregadas). \textbf{De la 917 en adelante no hay archivo.} Veintitrés actos del departamento; \textbf{ningún balance municipal} y \textbf{ningún plano}. \\
2 ene.\ 1922* & \textbf{Decreto 176}: declara cesante al subcomisario de Mojotoro \textbf{Medardo Sarmiento} «\textit{en vista del continuo abandono que hace de su puesto}» y nombra a \textbf{Rodolfo Pérez}. \textbf{Es el único relevo del tramo cuyo motivo el Boletín publica}, y el acto no nombra al departamento. \\
20 ene.\ 1922* & La Junta Electoral publica los \textbf{circuitos del Colegio Electoral Nº 3}: cuatro mesas, los mismos cuatro edificios desde 1918 y \textbf{treinta y cuatro parajes}. Donde 1918 imprime \textbf{Calderilla}, 1922 imprime \textbf{Candelaria}; se transcribe y no se concilia. \\
8 mar.\ 1922* & \textbf{Decreto 349}: declara cesante al comisario de policía del departamento \textbf{Julio Royo} y nombra a \textbf{Jorge A.\ Rauch}. Tercer comisario en tres meses y medio, y ninguno de los tres actos da la causa. \\
21 y 22 mar.\ 1922* & \textbf{Decreto 399}: subcomisario de Mojotoro \textbf{Juan de la Cruz Nogales} y subcomisario \textit{ad honorem} de \textbf{Vaqueros}, \textbf{Eustaquio Guerra}. \textbf{Decreto 402}: cesante \textbf{Casimiro Cáceres} del Registro Civil del departamento; entra \textbf{Julio Royo Ortiz}. \\
18 abr.\ 1922* & \textbf{Decreto 510}: el Interventor Nacional designa las Comisiones Municipales de los departamentos «\textit{que en la actualidad carecen de ella}». La de La Caldera va primera: \textbf{Daniel Linares\index{Linares, Daniel}, Augusto Regis\index{Regis, Augusto}, Domingo Royo, Jorge Rauch y Manuel Condori}. \textbf{El departamento no tenía órgano municipal antes de esa fecha, y el acto lo dice.} \\
20 abr.\ 1922* & \textbf{Decreto 524}: deja sin efecto el nombramiento de «D.\ Julio Royo» en el Registro Civil \textbf{por no haber tomado posesión} y nombra a \textbf{Santiago Guzmán} (expte.\ 4054-E). Ningún acto del tramo dice después qué pasó con Guzmán. \\
28 may.\ 1922* & La Comisión Municipal recién creada \textbf{sesiona y sanciona el presupuesto de gastos y cálculo de recursos} del departamento. \\
jun.\ 1922* & Un decreto del gobernador \textbf{Adolfo Güemes} \textbf{rectifica el artículo 2º del decreto 510} y repite los mismos cinco nombres, con otro orden y una inicial de más. \textbf{Qué rectifica no consta en el acto.} \\
12 jun.\ 1922* & \textbf{Decreto 105}: «\textit{vista la terna elevada por la Comisión Municipal de La Caldera}», jueces de paz \textbf{Carlos F.\ Matienzo}\index{Matienzo, Carlos} (propietario) y \textbf{Vicente Vivas} (suplente). \textbf{Cierra una vacante de seis meses} abierta por la renuncia de Laureano Peralta el 17 de diciembre de 1921 (decreto 123), y es el primer acto del tramo en que un órgano del departamento propone y no sólo recibe. \\
27 jun.\ 1922* & \textbf{Decreto 132}: acepta la renuncia de \textbf{Augusto Regis}\index{Regis, Augusto} a la Comisión Municipal, «\textit{atento los motivos que la determinan}», que no transcribe. \textbf{Dos meses y nueve días}, y el tramo termina sin acto que cubra la vacante. \textbf{Decreto 135}: \textbf{Carlos F.\ Matienzo}\index{Matienzo, Carlos} es nombrado encargado del Registro Civil del departamento, quince días después de ser su juez de paz propietario. \textbf{Cuarto encargado del Registro Civil en cuatro meses.} \\
15 jul.\ 1922* & \textbf{Decreto 172}: el Poder Ejecutivo \textbf{aprueba el presupuesto} que la Comisión Municipal sancionó el 28 de mayo, previo dictamen del Fiscal General. \textbf{El presupuesto aprobado no se publica.} \\
28 jul.\ 1922 & Fecha fijada para el \textbf{remate de la finca Potrero del Gallinato}, base \$800, en la ejecución Alsina c/ José Nogales (aviso Nº 264, B.O.\ Nº 910, h.\ 16). \textbf{Cae después del final del tramo leído y el archivo no dice si se vendió.} \\
1923 & \textbf{Año leído edición por edición} (52 ediciones, 939 a 990, 673 hojas; una por cada viernes del año, \textbf{sin un solo hueco en la serie numérica}). Diecisiete actos del departamento; \textbf{ningún plano} y \textbf{ningún balance municipal}. El original imprime el nombre del departamento con las letras separadas, y tres actos ---uno de ellos el único de dominio--- sólo aparecen con una búsqueda que admite espacios entre todas las letras. \\
18 ene.\ 1923* & \textbf{Decreto 620}: acepta la renuncia de \textbf{Augusto Regis}\index{Regis, Augusto} a la Comisión Municipal ---la segunda en siete meses--- y designa a \textbf{Genaro Molina}, comisario auxiliar de la Calderilla entre 1909 y 1913. El decreto imprime el año «1823» [\textit{sic}]. \\
22 ene.\ 1923* & \textbf{Decreto 629}: convoca a la renovación de la Legislatura para el 4 de marzo. \textbf{Caldera elige un senador y un diputado.} \\
9 feb.\ 1923* & La Junta Electoral publica los \textbf{circuitos del Colegio Electoral Nº 3}: cuatro mesas, los cuatro edificios de siempre y \textbf{treinta y cuatro parajes}. \textbf{Calderilla vuelve al circuito de la mesa 2}, donde 1922 imprimía «Candelaria». \\
17 feb.\ 1923* & Dos partidos proclaman por el departamento: la U.C.R.\ Nacional al senador Ignacio Ortiz y al diputado Augusto P.\ Matienzo; la Unión Provincial, al senador \textbf{José Antonio Aráoz} ---su tercer intento con el mismo candidato desde 1920--- y al diputado Pedro J.\ Frías. \\
21 feb.\ a 1 mar.\ 1923* & \textbf{Tres comisarios de policía del departamento en nueve días}: renuncia de \textbf{Eustaquio Murúa}\index{Murúa, Eustaquio} (decreto 701), \textbf{Juan de la Cruz Nogales} interino (712) y, tres días después, revocación de ese nombramiento y designación de \textbf{Julio Royo}, comisario de La Silleta (725). El decreto que revoca fecha el revocado el 24 de febrero y el revocado se fecha el 26. \\
may.\ 1923* & El Senado incorpora por Caldera a \textbf{José A.\ Aráoz}. \textbf{Del diputado electo no hay noticia} en las 673 hojas del año. \\
7 may.\ 1923* & \textbf{Decreto 886}: acepta la renuncia de \textbf{Domingo Royo} a la Comisión Municipal y nombra a \textbf{Julio Royo Ortiz}. Segunda vacante del año en ese cuerpo. \\
8 ago.\ 1923* & \textbf{Decreto 1076}: el Ejército clausura el camino que cruza el «Campo General Belgrano» hacia la \textbf{Quebrada de Lesser} y se abre otro costeando el \textbf{río Vaqueros}, «por haber sido éste siempre el camino nacional». \textbf{El acto no nombra al departamento.} \\
24 sep.\ 1923* & \textbf{Decreto 1161}: \textbf{declara intervenida la Municipalidad del departamento de La Caldera} y nombra Interventor a \textbf{Eduardo F.\ Harbin}, a pedido de presentaciones que instruye el expediente 4541 F. \textbf{Ni el expediente, ni las presentaciones, ni las instrucciones al Interventor, ni un solo acto suyo se publican.} Diecisiete meses después de que un Interventor Nacional creara ese cuerpo porque el departamento carecía de él. \\
2 nov.\ 1923* & Aviso de remate de la \textbf{finca Calderilla}, base \$7.000, en el juicio sucesorio de \textbf{Silvano Murúa}\index{Murúa, Silvano Ignacio}; con ella, la novena parte de una casa de la capital y 217 vacunos. Martillero Leguizamón; subasta fijada para el 11 de diciembre. \textbf{Primer acto de dominio del departamento desde 1921.} \\
28 dic.\ 1923* & Se publica por edicto el \textbf{pedido de concesión de agua de Francisco S.\ Urquiza}\index{Urquiza, Francisco S.} sobre los sobrantes del río Wierna para «Cerro de Buena Vista»: \textbf{cien litros por segundo para ciento noventa hectáreas}, con croquis que no se publica y una acequia que atraviesa la finca Vaqueros «de propiedad de los herederos Toranzos y doctor \textbf{Carlos Serrey}»\index{Serrey, Carlos}. El escrito es del 10 de diciembre de 1922 y la providencia, del 12 de enero de 1923, mandó publicarlo \textbf{en el diario \textit{La Provincia}}. \\
1924 & \textbf{Año leído edición por edición} (52 ediciones, 991 a 1042, 792 hojas): \textbf{el primer año de toda la serie en que el calendario cierra solo}, un viernes por edición y sin huecos. \textbf{Veinticuatro actos} del departamento en veinte ediciones ---menos que 1918, con veintisiete, y que 1925, con veinticinco---; \textbf{ningún plano} y \textbf{ningún balance municipal}. \\
1924 & Presupuesto General de la Provincia: el encargado del Registro Civil del departamento, con el sueldo común de las oficinas de campaña, y \textbf{mil pesos para la Iglesia de La Caldera}. \textbf{Es la única cifra propia que el Estado provincial le asigna al departamento por su nombre en todo el año.} \\
7 y 14 ene.\ 1924* & \textbf{Decreto 1381}: renuncia de \textbf{Julio Royo Ortiz} a la Comisión Municipal. \textbf{Decreto 1397}, seis días después: el mismo hombre, nombrado ahora «Julio Royo», es \textbf{expendedor de guías y marcas} del departamento, con fianza del art.\ 77 de la Ley de Contabilidad. \\
28 ene.\ 1924* & \textbf{Decreto 1429} (expte.\ 3898 C): \textbf{concede a Francisco S.\ Urquiza}\index{Urquiza, Francisco S.} los sobrantes del río Wierna, cien litros por segundo, \textbf{precaria por no haberse hecho aforos}, con el fundamento de que el río, después de proveer a las fincas de Wierna y Vaqueros, deja correr sobrante abundante. \textbf{El informe municipal que la funda se emitió con la comuna intervenida.} \\
15 mar.\ 1924 & \textbf{Edicto sucesorio de don Simón Yufra}, dictado \textbf{en La Caldera} por el juez de paz \textbf{Carlos F.\ Matienzo}\index{Matienzo, Carlos}. \textbf{Único acto de todo el año dictado dentro del departamento y fechado allí.} El apellido Yufra vuelve en el revalúo de 1944. \\
21 mar.\ 1924* & Edicto de \textbf{deslinde de la finca «Vaqueros» o «Entre Ríos»} (auto del 5 de diciembre de 1923), pedido por Natal López Cross: al este, propiedad de \textbf{Francisco Urquiza}\index{Urquiza, Francisco S.}; al norte, los herederos de Anacleto Toranzos y Teodolina Torino de Toranzos; al oeste, el camino nacional. \\
3 y 15 abr.\ 1924* & \textbf{Decreto 1560}: renuncias «indeclinables» de \textbf{Daniel Linares}\index{Linares, Daniel} y \textbf{Genaro Molina} a la Comisión Municipal. \textbf{Decreto 1576} (expte.\ 4969 F), doce días después: la integra con \textbf{Julio Royo Ortiz} ---que había renunciado en enero---, \textbf{Camilo Cáceres}, \textbf{Francisco Urquiza}\index{Urquiza, Francisco S.} y \textbf{Serafín Yáñez}. \textbf{Tres meses después de que se le concediera el agua, el concesionario integra el cuerpo cuyo informe fundó la concesión.} \\
16 abr.\ 1924* & \textbf{Decreto 1580} (expte.\ 3247 D): aquiescencia para instalar una \textbf{escuela de la Ley 4874 (Láinez) en «San Alejo»}. \textbf{Primera escuela nacional del departamento} que el relevamiento encuentra. \\
1 jul.\ 1924* & \textbf{Resolución del Ministro de Gobierno}, sin número: autoriza a \textbf{investigar el destino que hayan tenido los libros del Registro Civil} del departamento, por urgente y por los «intereses que el mismo puede comprometer», con cooperación de la comisaría. \textbf{No se publica el resultado.} \\
31 jul.\ y 4 ago.\ 1924* & \textbf{Decreto 1786}: jueces de paz \textbf{Carlos F.\ Matienzo}\index{Matienzo, Carlos} y \textbf{Vicente Vivas}, por terna de la Comisión Municipal ---el mismo par de 1922, sin que el decreto cite al anterior ni a la intervención del medio---. \textbf{Decreto 1796} (expte.\ 5180 F): renuncia de \textbf{Serafín Yáñez}, a los cuatro meses de ser designado, \textbf{sin reemplazante publicado}. \\
14 nov.\ 1924* & \textbf{Decreto 1975} (expte.\ 5162 F): \textbf{aprueba el presupuesto municipal} de La Caldera del ejercicio en curso, con dictamen del Fiscal General. \textbf{No publica una sola cifra}, y lo aprueba a seis semanas del cierre del ejercicio. \\
1925 & \textbf{Año leído edición por edición} (52 ediciones, 1043 a 1094, 798 hojas, sin huecos). \textbf{Veinticinco actos} del departamento; \textbf{ningún plano} y \textbf{ningún balance municipal}. \\
1925 & Al menos seis pedimentos de cateo en enero, tres de ellos publicados con el propietario del suelo; ocho solicitudes del año caducarán, y la de Larrán prosperará. Se disuelve Levilly y Cía.; se remata la fundición con base de \$222.144. \\
ene.--feb.\ 1925* & \textbf{El departamento se parte en dos colegios electorales}: \textbf{Caldera, Colegio Nº 5} (Escuela Elemental Mixta e Iglesia Parroquial) y \textbf{Vaqueros, Colegio Nº 6} (Escuela Infantil Provincial y Oficina de Correos de Mojotoro). De 1918 a 1924 había sido siempre el Colegio Nº 3, con cuatro mesas y un solo rótulo. \textbf{Y desaparece la nómina de parajes}: los actos fijan edificios y autoridades y ya no enumeran circuitos. Caldera elige \textbf{un diputado}. \\
11 y 12 may.\ 1925* & \textbf{Decreto 2453}: \textbf{Carlos Murúa}\index{Murúa, Carlos} comisario de policía del departamento, en reemplazo de Julio Royo Ortiz. \textbf{Decreto 2465}, al día siguiente: el mismo Murúa, \textbf{Receptor de Rentas y Recaudador de Impuestos al Consumo}, con fianza de \$5.000. \textbf{Tercera vez en nueve años que una sola persona junta la comisaría y la caja del departamento.} \\
2 y 25 jun.\ 1925* & \textbf{Decreto 2575}: renuncia de \textbf{Julio Royo Ortiz} a la Comisión Municipal. \textbf{Decreto 2671}: «encontrándose \textbf{desintegrada}» la Comisión, la integra con \textbf{Daniel Linares}\index{Linares, Daniel}, \textbf{Rafael Delgado} y \textbf{Genaro Molina} ---Linares y Molina son los dos que habían renunciado «indeclinablemente» catorce meses antes---. \\
25 jun.\ 1925* & \textbf{Decreto 2674}: subcomisario \textit{ad honorem} de Vaqueros, Antonio Agüero (hijo), \textbf{en reemplazo de Francisco S.\ Urquiza}\index{Urquiza, Francisco S.}. \textbf{Quinta calidad del concesionario del agua en año y medio}, y el acto que lo saca no dice desde cuándo estaba: \textbf{su nombramiento no aparece en ningún tramo leído.} \\
1 ago.\ 1925* & \textbf{Auto de concesión del cateo 1168-C a Juan Larrán} sobre la finca Las Nieves, de Garzón y Pintos, y terrenos de \textbf{E.\ Quipildor}; cuatro unidades, 2.000 ha, anotado bajo el Nº 87. \textbf{Es el único de los nueve pedimentos de enero de 1925 cuyo auto de concesión consta.} Los edictos citan un deslinde de \textbf{1889} del agrimensor P.\ Arias: el documento más antiguo que el libro ve citado sobre el departamento. \\
14 ago.\ 1925* & Edicto de deslinde de «Sauce» o «Paraíso» y la fracción \textbf{Ceival}, del \textbf{Dr.\ José María Solá}, parte en Campo Santo y parte en La Caldera. Linderos: «Mosquera», \textbf{de Daniel Linares}\index{Linares, Daniel} ---miembro de la Comisión Municipal desde hacía dos meses---, los «Noques» de Reymundo Echenique y \textbf{«La Despensa», de los herederos Murúa}. \\
29 oct.\ 1925* & \textbf{Decreto 2962}: renuncia de Carlos Murúa\index{Murúa, Carlos} a la comisaría, «dándosele las gracias por los servicios prestados»; entra Francisco Arias. \textbf{El decreto no dice nada de la receptoría de rentas.} \\
1926 & \textbf{Año leído edición por edición} (52 ediciones, 1096 a 1147, 769 hojas, sin huecos). \textbf{Dieciocho actos} del departamento. Los cinco resúmenes de Tesorería del año no son meses consecutivos y \textbf{no permiten armar ninguna cadena}. \\
1926 & Otros actos del año: refacción y alquiler de la casa de la comisaría (decretos 3205 y 3524); jueces de paz \textbf{Matienzo} y Molina (3248) y renuncia del suplente Lucas N.\ Molina (3322); tres encargados del Registro Civil en cinco meses (3252, 3350 y 3504); sucesorio de Celestina Velázquez de Ruiloba ante el juez de paz; citación por \textbf{reposición del título de «Los Peñones»}; deslinde de una finca sin nombre lindera al río Wierna. \\
1926 & Juicio para \textbf{reponer el título perdido} de la finca ``Peñones o Los Peñones''. \\
30 ene.\ 1926* & \textbf{Decreto 3156}: vuelve la nómina de parajes y vuelve \textbf{partida}. «La Caldera, Colegio Electoral Nº 5», dos mesas y \textbf{veintiséis parajes}; a continuación, en el mismo acto, «Vaqueros, Colegio Electoral Nº 6». Los ocho nombres que faltan respecto de 1924 ---Lesser, Vaqueros, Mojotoro, Mojotorito, Potrero de Gallinato, Túnel y Estación--- pasaron al Colegio 6, y «El Monte de las Garrapatas» se imprime como un solo topónimo. \\
4 may.\ 1926* & \textbf{Decreto 3420}: aprueba el presupuesto de la Comisión Municipal de La Caldera. \textbf{Tercer presupuesto municipal aprobado sin publicar una cifra}, después de 1922 y 1924. \\
1927 & \textbf{Año leído edición por edición} (52 ediciones, 1148 a 1199, 745 hojas, sin huecos). \textbf{Veinticuatro actos} del departamento. La cadena de saldos se arma en dos tramos: faltan junio, julio, agosto y diciembre. \\
1927 & Remate del Banco de la Nación: \textbf{una casa quinta de diez hectáreas en Lesser}, base \$6.000. La segunda residencia caldereña está documentada desde entonces. \\
1927 & \textbf{Bernardo Austerlitz\index{Austerlitz, Bernardo}}, propietario de la finca San Alejo y Santa Rufina según los pedimentos de cateo de 1925, es electo \textbf{senador por el departamento de Caldera}. \\
11 feb.\ 1927* & Se republica la nómina del \textbf{Colegio Electoral Nº 5}: los mismos veintiséis parajes de 1926. \\
17 y 20 may.\ 1927* & \textbf{Decreto 5024} (expte.\ 1099 M): aprueba la \textbf{rendición de cuentas de la Comisión Municipal de La Caldera por el ejercicio de 1926}, atento el informe del Consejo General de Educación «en lo que respecta a la \textbf{renta escolar} que por dicho ejercicio le corresponde». \textbf{Decreto 5036} (expte.\ 1725 M): aprueba su presupuesto de 1927. \textbf{Los dos aprueban sin publicar una sola cifra}, y el segundo sale en una tanda con General Güemes y La Poma. \\
jun.--ago.\ 1927 & La finca \textbf{«Los Sauces»} se remata en el \textbf{juicio sucesorio de Domingo Royo} ---miembro de la Comisión Municipal entre 1922 y 1923---: base \textbf{\$50.000} el 3 de junio, \textbf{\$14.000} el 8 de julio y el mismo aviso otra vez el 26 de agosto, con tres números de aviso distintos. \textbf{No consta el resultado.} \\
13 jun.\ 1927 & Edicto sucesorio de \textbf{Milagro Catacata}, dictado y firmado \textbf{en La Caldera} por el juez de paz \textbf{Ángel Fernández}. \\
1928 & \textbf{Año leído edición por edición} (52 ediciones, 1200 a 1251). \textbf{Veintitrés actos} del departamento. Faltan los resúmenes de Tesorería de marzo y octubre; \textbf{la serie de decretos salta del 7097 al 8001 en un solo día}. \\
1928 & En tres semanas: se rechaza el reclamo de Campo Santo contra la concesión de agua de Urquiza\index{Urquiza, Francisco S.}; \textbf{Urquiza renuncia a la Comisión Municipal de La Caldera}; la comisión es declarada en acefalía; y un nuevo gobernador \textbf{revoca la concesión} por ser ``inconveniente''. \\
1928 & Agosto: Urquiza\index{Urquiza, Francisco S.} es sometido a \textbf{proceso penal por ``hurto y uso indebido de agua'', con prisión preventiva}, por denuncia de \textbf{Juan A. Bianchi\index{Bianchi, Juan A.}}, dueño de la finca vecina. Sobreseído por prescripción en 1930. \\
26 mar.\ 1928* & \textbf{Decretos 6776 y 6777}: renuncian el \textbf{juez de paz propietario Victorio Offredi} y el \textbf{comisario Ramón Balvoa}, sin causa declarada. \textbf{Seis semanas antes de la acefalía.} \\
9 may.\ 1928* & \textbf{Decreto 7096}: «debiendo procederse a la reorganización de la Municipalidad de La Caldera, \textbf{para asegurar la existencia del régimen municipal}», declara \textbf{en acefalía} la Comisión Municipal y nombra comisionado \textit{ad honorem} a \textbf{Luis E.\ Guardo}. El decreto anterior, 7094, hace lo mismo con Embarcación. \textbf{Decreto 7095}: comisario, \textbf{Julio Royo Ortiz}. \textbf{Decreto 8002}: Registro Civil, \textbf{Teófilo Reyes} por Ramón Balboa. \\
9 y 12 may.\ 1928* & \textbf{Decreto 8001}: subcomisario \textit{ad honorem} de Mojotoro, \textbf{Eustaquio Guerra}. \textbf{Decreto 8024}, tres días después: acepta su renuncia y nombra a Juan de la Cruz Nogales. \textbf{Nombramiento y renuncia salen en la misma edición, en hojas contiguas y sin remisión entre sí.} \\
1 ago.\ 1928* & \textbf{Decreto 9464}: reconoce servicios a los ex-encargados del Registro Civil; el renglón de La Caldera le asigna a Ramón Balboa un período que \textbf{empieza el 9 de junio y termina el 13 de mayo}. Se transcribe y no se concilia. \\
19 oct.\ 1928* & \textbf{Decreto 9750}: juez de paz suplente, \textbf{Teófilo Reyes}, por terna de la Comisión Municipal ---el mismo que en mayo había tomado el Registro Civil---. \textbf{Siete cargos del departamento renovados en siete meses, y ninguno queda en pie.} \\
1929--1932 & \textbf{Cuatro años leídos edición por edición} (ediciones 1252 a 1461). \textbf{Ninguno de los cuatro publica un solo plano.} El departamento pasa al \textbf{Colegio Electoral Nº 2}, el tercero en once años. \\
1929--1934 & \textbf{Tres mujeres en dos oficinas del departamento}, donde el archivo temprano no registra ninguna entre 1908 y 1928: \textbf{Clara Iradi} en la receptoría de rentas (1929 y 1931), \textbf{Julia Regis de Regis} en la misma receptoría (1933) y \textbf{Laura Aráoz de Murúa} en el Registro Civil (1934). \\
1929 & \textbf{Adeodato Aybar\index{Aybar, Adeodato} termina la catastración} de los departamentos de Campo Santo, Caldera y La Viña. Ley 1269 (original 11078): se autoriza a la Municipalidad de \textbf{Campo Santo} a intervenir el agua desde las juntas del Wierna con el de La Caldera, vigilar y prohibir nuevas tomas. \\
1929 & Remate de una \textbf{segunda casa quinta}, de cinco hectáreas con frente sobre el río Vaqueros. \\
22 mar.\ 1929* & \textbf{Decreto 10.389}: comisario de policía del departamento, \textbf{Manuel Condori}. \textbf{Primer apellido de origen andino al frente de la comisaría}, y uno de los seis que el revalúo de 1944 registrará entre sus titulares. \\
27 abr.\ 1929* & \textbf{Decreto 10.514}: receptoría de rentas del departamento, \textbf{Clara Iradi}. \textbf{Primera mujer con cargo público en La Caldera} que el archivo temprano registra, y el cargo es el de la caja. \\
5 ago.\ 1929* & \textbf{Decreto 10.910} (expte.\ 3463 D): liquida \$500 a cuenta al catastrador \textbf{Adeodato Aybar}, que declara \textbf{haber terminado la catastración de Campo Santo, Caldera y La Viña}, con imputación a la Ley 9720 de 1928. \textbf{Primer catastro del departamento que el libro encuentra}, del que no se publica superficie, parcelas ni plano, y cuya remuneración queda sin fijar. \\
27 sep.\ 1929* & \textbf{Ley 11.078}: autoriza a la \textbf{Municipalidad de Campo Santo} a trabajar la playa del río Mojotoro «para reconcentrar el agua, \textbf{desde las juntas del Río de Wierna con el de La Caldera}», y a \textbf{vigilar y prohibir nuevas bocas-tomas}. \textbf{El punto de arranque está dentro de La Caldera y la facultad es de otro municipio.} \\
27 sep.\ 1929* & \textbf{Ley 11.088}: \$500 para comprar un terreno en el \textbf{pueblo de La Caldera, «Capital del mismo Departamento»}, con destino a la \textbf{Casa Municipal}. \textbf{Veintiún años después del presupuesto de 1908, el municipio no tenía sede.} El archivo no dice si el terreno se compró. \\
27 sep.\ 1929* & \textbf{Ley 11088}: autoriza invertir \textbf{\$500} para adquirir un terreno en el pueblo de La Caldera destinado a la \textbf{Casa Municipal}. Lo mismo que seis años después se presupuestará para el espigón de emergencia contra el río, que terminó costando \$700. \\
1930 & Segundo remate de la \textbf{estancia Potrero de Gallinato}, junto a una casa de la capital, en la ejecución de un mismo deudor. Juan Larrán manifiesta el descubrimiento de las minas de \textbf{plomo y plata} Caldera Nº 1, 2 y 3 sobre la finca San José. \\
1930 & La Dirección de Minas declara caducas por abandono \textbf{las ocho solicitudes de cateo de 1925}. Ninguna llegó a explorar. \\
28 abr.\ 1930* & \textbf{Resolución 434} del Ministerio de Gobierno: declara \textbf{procedente el impuesto de degolladura} que la Municipalidad de La Caldera aplica a los abastecedores de las cuadrillas del camino nacional Salta--Jujuy dentro de su jurisdicción. Es el único acto de las ediciones leídas en que el municipio cobra un tributo a un tercero y la Provincia se lo confirma; de sus ordenanzas de impuestos el Boletín sólo publica las aprobaciones, sin articulado. Llega por una resistencia al pago. \\
sep.--dic.\ 1930* & \textbf{Segunda intervención federal sobre la provincia en ocho años}: desde octubre los actos del departamento ---Comisión Municipal, subcomisarías, nombramientos--- los firma el \textbf{Interventor Nacional}. \\
oct.\ 1930* & La Dirección General de Minas concede a \textbf{Alberto Jaime Harrison y Juan Larrán} un cateo de 2.000 ha en la finca \textbf{Las Nieves}, en terrenos de Julio Suárez y Francisco de la Vega (expte.\ 38-H, B.O.\ Nº 1.346). \\
1931 & \textbf{Decreto 14.205} del Interventor Nacional: fija la categoría de los municipios de la provincia. \textbf{Caldera queda en tercera}, gobernada por comisión con presidente designado por el Poder Ejecutivo y mandato de un año. Sigue clasificada así en 2022. \\
1931 & La \textbf{Corte de Justicia} deja sin efecto la revocación de 1928 y ordena restablecer la concesión de Urquiza\index{Urquiza, Francisco S.}. Doctrina: los motivos de interés público ``deben aparecer establecidos por formalidad y no meramente invocados''. \\
mar.\ 1931* & La Dirección General de Obras Públicas «ratifica la necesidad de confeccionar el \textbf{plano catastral de la Provincia} [\ldots] en láminas que comprendiendo una cada departamento». \textbf{Citado y no acompañado}, y dos años después de la catastración que Aybar dio por terminada. \\
23 oct.\ 1931* & \textbf{Decreto 14091}: nombra una Comisión ad-hoc ---el Director General de Obras Públicas y los Comisionados Municipales de Campo Santo y de La Caldera, \textbf{Augusto Regis}\index{Regis, Augusto}--- para \textbf{investigar todas las concesiones de agua de los ríos Vaqueros y Wierna}, las ordenanzas municipales que las reglaban y los usos y costumbres. Sus considerandos declaran que la distribución del agua de riego es atribución municipal «mientras no se sancione una ley general de Irrigación». \textbf{Este libro no encontró el informe.} \\
30 oct.\ 1931* & Decreto que crea una \textbf{Comisión \textit{ad hoc} sobre las aguas de los ríos Vaqueros y Wierna}, afluentes del Mojotoro, \textbf{integrada por los Comisionados Municipales de Campo Santo y de La Caldera} ---los dos designados, ninguno electo---. \\
1932 & Deslinde judicial de cinco fincas pedido por Silvano Ignacio Murúa y su esposa, Fanny Ovejero, entre ellas \textbf{``El Durazno o Peñones''}. Los linderos nombrados incluyen a \textbf{Daniel Linares\index{Linares, Daniel}} --- presidente de la comisión de defensas de 1910 --- como dueño de Los Porongos. \\
1932 & Decreto 15.445: el P.E. aprueba \textbf{en octubre} el presupuesto municipal ``para regir durante el ejercicio 1932'', y ordena tramitar el saldo deudor por \textbf{renta escolar atrasada}. \\
1933 & Vialidad de Salta encomienda a \textbf{Augusto Regis\index{Regis, Augusto}}, presidente del municipio, la limpieza del camino de La Caldera a Tres Cruces por el bajo, con quinientos pesos y rendición documentada: \textbf{el municipio como contratista vial de otra jurisdicción}, ochenta y un años antes del contrato de 2014. \\
1933 & \textbf{Augusto Regis\index{Regis, Augusto}} --- jurado en 1908, comisionado para las defensas del río en 1910, miembro de la comisión municipal en 1922 --- es renombrado \textbf{Presidente de la Comisión Municipal}. El padrón minero registra las minas Caldera 1, 2 y 3 como ``a mensurarse''. \\
oct.\ 1933 & Decreto 16.798, del 2 de octubre, publicado tres meses después: se aprueba el presupuesto municipal de 1933 pese a que \textbf{contraría el art. 79 de la Ley 68}, que prohíbe destinar más del 25\,\% de la renta a sueldos administrativos. \\
27 oct.\ 1933* & Alberto, Abel y Adolfo Murúa ponen en remate \textbf{por su cuenta} la finca \textbf{La Despensa}, 2.700 hectáreas con riego, base \$35.000 sobre una avaluación fiscal de \$50.000 (B.O.\ Nº 1.503). El aviso nombra como vecinos a la sucesión y a los herederos de Silvano Murúa, padre de Silvano Ignacio. \\
1934 & \textbf{Ley 1423 (orig.\ 145)}: La Caldera integra la ``región económica del olivo'' y se eximen \textbf{por diez años de todo impuesto provincial o municipal} las tierras donde se planten olivos. Deslinde del Puesto Despensa: las cinco fincas de 1932 ya están \textbf{adjudicadas entre herederos}. \\
1934 & En cinco días, dos decretos más clasifican al departamento en el último escalón: el Registro Civil en \textbf{tercera de tres} categorías y la Receptoría de Rentas en \textbf{quinta de cinco}. \\
1935 & Nuevas defensas: se liquidan \$500 al presidente municipal para un \textbf{espigón de emergencia ``pie de gallo''}. El decreto consigna que \textbf{la solución definitiva sería un espigón macizo} que ``alejaría por completo el peligro'' y que debe estudiarse. Un mes después la partida se amplía en \$200 más: \textbf{cuarenta por ciento de sobrecosto}, el mismo mecanismo que en 1911, cuando fue del cuarenta y ocho. \\
1935 & La Dirección General de Obras Públicas licita la \textbf{adquisición de terrenos} para edificación escolar --- entre ellos en \textbf{Vaqueros} --- y para asistencia social --- entre ellos en \textbf{La Caldera}. \\
1935 & Remate por ejecución hipotecaria del Banco Provincial de la finca \textbf{``Vaqueros o Entre Ríos''}, base \$12.000. Linda al este con Francisco Urquiza\index{Urquiza, Francisco S.} y al oeste con el camino nacional. \\
1935 & Los \textbf{vecinos de La Caldera piden a Vialidad un puente} que una el pueblo con el camino nacional. Respuesta: se tomó nota, se dispondrá ``en su oportunidad''. No se construye. \\
1935 & Diciembre: Obras Públicas eleva un \textbf{proyecto completo con levantamiento, planos y presupuesto}. Identifica que el río se bifurca en tres brazos cuatro kilómetros aguas arriba y propone intervenir \textbf{en la bifurcación}, no junto a las casas. \$3.757,50 por administración. Declara que servirá para ``adquirir experiencia sobre el comportamiento del Río'' y proyectar después ``obras más duraderas''. \\
7 jun.\ 1935* & \textbf{Ley Nº 172}: autoriza a invertir \textbf{\$7.000 en reformar y ampliar la Iglesia Parroquial} del departamento, según planos del ing.\ Antonio Mille, con cargo a Rentas Generales. \textbf{Diez veces} lo liquidado ese mismo año para impedir que el río entrara al pueblo ---\$700 entre la partida inicial y su ampliación---, y casi el doble del proyecto completo de defensas que Obras Públicas elevó en diciembre. \\
14 jun.\ 1935 & \textbf{Inventario provincial de libros del Registro Civil}, por oficina. El renglón del departamento dice entero: «\textit{C[a]ld[e]ra.. matrimonios 1913 1926 13 400 116}». \textbf{Un solo libro, y de matrimonios}, cuando las otras oficinas traen dos o tres; el cuadro no dice por qué ni declara que falte ninguno. Once años después de la diligencia de 1924 sobre «el destino que hayan tenido los libros». \\
15 nov.\ 1935* & La Provincia transfiere a la Caja Nacional \textbf{\$2.500} de indemnización por el accidente del que fue víctima el obrero \textbf{José Gallardo} en el \textbf{camino Gallinato} (Ley 9.688). En marzo de 1936 el Directorio de Vialidad paga tres ataúdes por «\textit{los tres obreros fallecidos en el accidente de Gallinato}». Son las únicas muertes obreras del archivo. \\
1936 & Un decreto autoriza \$32,87 para \textbf{terminar de pagar} las defensas del pueblo sobre el río: la obra se hizo. \\
1936 & Remate \textbf{sin base} de un \textbf{aserradero completo} —sierra sin fin, dos circulares, motor a vapor, garlopa y dieciséis bueyes— en el pueblo de La Caldera, por deuda con una casa comercial de la capital. \\
1936 & Remate de la finca \textbf{La Despensa} en la ejecución de Salvador Rosa contra tres herederos Murúa, con los otros herederos como linderos. \textbf{Deslinde en 1932, intento de venta en 1933, partición en 1934, ejecución en 1936}: el ciclo completo por el cual un dominio grande se fragmenta. \\
1936 & La Corte \textbf{rechaza la demanda de daños} de Urquiza\index{Urquiza, Francisco S.}: aun ilegal la revocación, no probó que hubiera podido usar esa agua. \textbf{La misma ausencia de aforos que hizo precaria la concesión le impidió cobrar los daños.} \\
6 mar.\ 1936* & El Tribunal Electoral \textbf{anula la elección del 1º de marzo en la Mesa Nº 1 del Circuito 5 de La Caldera} ---la de la Escuela Elemental Mixta, presidida por \textbf{Augusto Regis}\index{Regis, Augusto}, con \textbf{Daniel Linares}\index{Linares, Daniel} como suplente segundo, ambos insaculados--- y se convoca a comicios complementarios para el 15 de marzo. El decreto no publica el fundamento de la anulación. \\
1937 & \textbf{Año leído edición por edición} (53 ediciones, 1669 a 1721). \textbf{Deslinde de «La Helvecia» antes «Calderilla»}, lindando al norte con la propiedad de \textbf{Augusto Regis}\index{Regis, Augusto} ---presidente de la Comisión Municipal ese mismo año---, al naciente con los herederos de Pedro Valencia y al poniente con el río de La Caldera. Se pagan las \textbf{reparaciones de las líneas telefónicas del pueblo}, a cargo de \textbf{Rafael Delgado}, juez de paz propietario del distrito. Se crea una comisaría de 2.ª en \textbf{San Alejo y Santa Rufina} y se pide un \textbf{guarda sanitario con asiento en La Caldera}. Decenas de órdenes de pago de Vialidad por los caminos \textbf{Calderilla--El Desmonte por Gallinato}, \textbf{Calderilla--Campo Santo} y \textbf{Calderilla--Güemes}. \\
1937 & Deslinde de la finca \textbf{``La Helvecia'', antes ``Calderilla''}: sus linderos son \textbf{Augusto Regis\index{Regis, Augusto}} --- nombrado presidente municipal veinte días después ---, Genaro Molina\index{Molina, Genaro}, los herederos de Pedro Valencia y \textbf{el río La Caldera}. Lo solicita Hermann Pfister\index{Pfister, Hermann}, perito de los deslindes de 1932 y 1934. \\
1937 & Remate de la finca \textbf{Chalchanio} —partido de 1909— en la ejecución Terán c/ Juárez. El edicto cita \textbf{folio, asiento y libro} del registro: es el primero de las ediciones leídas que remite a la inscripción en lugar de describir por lomas y mojones. \\
1937--1942 & \textbf{Gimnasio} de la Ley 439 de 1937 para el pueblo: la Junta de Educación Física lo paga en enero de 1942, a \$1.500, y lo recibe Augusto Regis, presidente de la Comisión Municipal. Es el único de los actos materiales de 1929--1937 cuya llegada está probada (capítulo \ref{cap:presencia}). \\
10 sep.\ 1937* & \textbf{Decreto 1421}: niega la aquiescencia para instalar una escuela nacional de la Ley 4874 en \textbf{Potrero de Castilla} alegando que el Consejo General de Educación «instalará próximamente escuelas fiscales» allí. \\
4 oct.\ 1937* & \textbf{Decreto 1500}: \textbf{veinticuatro días después}, deroga ese artículo y concede la aquiescencia, porque la Inspección Nacional verificó el porcentaje de analfabetos y comprobó que \textbf{no existía escuela provincial en la localidad ni en lugares próximos}. Potrero de Castilla es el paraje donde la comunidad kolla Kondorwaira obtendrá personería en 2009. \\
1938 & \textbf{Año leído edición por edición} (52 ediciones, 1722 a 1773). \textbf{Decreto de defensas sobre la margen izquierda del río Caldera}, a pedido de vecinos, porque «\textit{todos los años en la época de lluvias las aguas del río Caldera amenazan las propiedades}»: obra «defensa río Caldera», \textbf{cincuenta metros}, con croquis que no se publica. \textbf{Ochenta y seis años antes del Plan de Manejo de 2024.} \\
1938 & Se pide la \textbf{segunda concesión de agua del departamento}: regadío de la finca \textbf{«Campo Alegre»}, con título al folio 374 asiento 475 del libro D de La Caldera, \textbf{con intervención de la Municipalidad de La Caldera} y dictamen del Fiscal de Gobierno. El archivo no dice si se concedió. Y un pedimento minero de \textbf{2.000 hectáreas} se ubica sobre las fincas Rosal y Potrero, \textbf{de la sucesión Quipildor}, entre Rosario de Lerma y Caldera ---el mismo apellido que en 1925 era dueño de los terrenos del cateo de Larrán---. \\
1938 & Elección: por el departamento, senador \textbf{Florentino M.\ Serrey}\index{Serrey, Carlos} por un partido y \textbf{Francisco Mercado} con \textbf{Teófilo Reyes} de suplente por el otro. \textbf{Comisión Municipal}: titulares \textbf{Eustaquio Murúa}\index{Murúa, Eustaquio} y Eugenio Muñoz, con \textbf{Genaro Molina} entre los suplentes ---veintidós y veintinueve años apareciendo en este archivo---. La nómina de circuitos del \textbf{Colegio Electoral Nº 2} trae ya Helvecia, Getsemaní, Abra de la Sierra y Cañada Ancha, que la lámina \ref{fig:parajesnom} atribuye sólo a 1940; el barrido la entrega cortada y se reclama entera. \\
1938 & El Consejo de Educación informa que la Comisión Municipal \textbf{no efectuó ingreso alguno al fondo propio de educación} y el P.E.\ \textbf{devuelve el balance de 1937 con observaciones}; un mes después lo aprueba y fija el porcentaje sobre lo recaudado. \textbf{Sin publicar una sola cifra.} Y un acto del mismo año dice que no ha cumplido «ninguna de las municipalidades de la campaña, \textbf{a excepción de la de La Caldera}, que ocasionalmente lo ha hecho». Se transcriben las dos y no se concilian. \\
1938 & Se aprueba el balance municipal de 1937 y se fija en \$283,51 su aporte educativo, el diez por ciento de la renta afectable: \textbf{la renta del municipio rondaba los \$2.835}, menos de la mitad de lo que valía una casa quinta del departamento. \\
1938 & \textbf{Isaac Fernández}, comerciante de la capital, pide concesión de agua del río Santa Rufina para regar doscientas hectáreas de la finca \textbf{``Campo Alegre''}, comprada en 1937 e inscripta al folio 374 del Libro D. Linda al sur con Daniel Linares\index{Linares, Daniel} y al oeste con la sucesión Austerlitz\index{Austerlitz, Bernardo}. \textbf{La finca está sin mensurar.} \\
1938 & Los \textbf{propietarios y vecinos de La Calderilla} piden defensas \textbf{sobre la margen izquierda} del río, ``en salvaguardia de las plantaciones''. Cincuenta metros, \$250, por administración. Es la margen hacia la que el informe de 1935 proponía desviar el caudal. \\
1939 & Al resolverse la concesión de «Campo Alegre», la oposición invoca «\textit{la existencia de una \textbf{ordenanza dictada por el municipio de la Caldera en fecha octubre 13 de 1914}, que reglamenta la distribución de la totalidad}» del caudal. \textbf{Veinticinco años después, la ordenanza de 1914 se esgrime como derecho vigente en un expediente provincial.} \\
1939 & \textbf{Augusto Regis}\index{Regis, Augusto} firma el certificado sobre la provisión de agua que el libro ya citaba, y es reelecto presidente de la Comisión Municipal. El comisario \textbf{Abraham Alegre} es destituido por «\textit{actos graves de inconducta}» y lo reemplaza \textbf{Claudio Murúa}\index{Murúa, Carlos}. Escuela de la Ley 4874 en \textbf{San Alejo}. \\
1939 & La \textbf{Curia Eclesiástica de Salta} promueve posesión treintañal sobre un terreno del pueblo de La Caldera. Sus linderos norte y oeste son el \textbf{Dr.\ Carlos Serrey\index{Serrey, Carlos}}. \\
1939 & Se \textbf{deniega} la concesión de agua para Campo Alegre: hubo disenso municipal y tres oposiciones particulares, y otorgarla alteraría el reparto fijado por la ordenanza de 1914. \\
1939 & Se rechaza inscribir un derecho de agua de la finca «Abra de Lesser»: el arroyo «nace y muere dentro de la finca» y por lo tanto es \textbf{de dominio privado}. \textbf{Había cursos enteros fuera de todo registro.} \\
1940 & \textbf{La defensa del río de 1938 no se hizo.} El decreto que autoriza «\textit{con carácter de urgencia}» una nueva dice que «\textit{el año ppdo.\ debieron efectuarse las obras de defensa en el río de la Caldera, \textbf{las que fueron suspendidas en virtud de que el mencionado río desvió su curso}}», y que «\textit{hoy nuevamente amenazan peligro las aguas de este río \textbf{para la población de la Calderilla}}». Autoriza un \textbf{reparo provisorio} en el lugar denominado la Calderilla. \textbf{El archivo no registra que se haya ejecutado.} \\
1940 & El cuerpo municipal pasa a llamarse de \textbf{concejales}: titulares Eugenio Muñoz y José Manuel Alfaro. \textbf{Augusto Regis}\index{Regis, Augusto} firma como «presidente de la comisión municipal del distrito de la Caldera, \textbf{diputado provincial}». Escuela de la Ley 4874 en \textbf{Gallinato}. En la mesa 1 del circuito 5 es suplente \textbf{Severino Mamani}: \textbf{primer apellido de origen andino en una autoridad de mesa del departamento}. La red vial provincial incorpora «\textit{de la Calderilla a río Saladillo por Gallinato y Campo Santo}» y la ruta que pasa por «\textit{Vaqueros, Calderilla, hasta Tres Cruces en el límite con Jujuy}». \\
1940 & Enero: se \textbf{confirma la denegatoria} de la concesión de Campo Alegre por «razones de orden legal que impiden al Poder Ejecutivo otorgar concesiones sobre el agua del Río Santa Rufina». Febrero: las defensas proyectadas en 1939 se habían suspendido \textbf{porque el río desvió su curso}; se autorizan \$500 para un reparo provisorio en La Calderilla. Mayo: el municipio debe \$951,95 a la Dirección Provincial de Sanidad. \\
30 ene.\ 1940* & Decreto de ubicación de mesas receptoras: La Caldera es el \textbf{Colegio Electoral Nº 2}, con dos circuitos, tres mesas y \textbf{cuarenta y seis parajes} nombrados uno por uno ---entre ellos \textbf{Getsemaní}, Las Lagunas, Mosquera, Potrero de Valencia, Potrero de Gallinato, Los Porongos, Campo Alegre, Lesser y Villa Urquiza\index{Urquiza, Francisco S.}---. Es la lista más completa de unidades territoriales pobladas entre 1909 y 2026. \\
18 jul.\ 1940* & La Provincia autoriza \$673,45 para traer al Jefe de la División de Piscicultura de la Nación y comprar diez mil truchas. Su informe nombra, entre las \textbf{aguas privadas} ya pobladas de pejerrey, la finca de \textbf{Lucio Ortiz en Yacones}, y su canal como una de las cuatro instalaciones disponibles de la provincia. \\
12 nov.\ 1940* & \textbf{Decreto 4217}: aprueba el Acta 376 de Vialidad de Salta con un \textbf{plan vial de quince años} (Ley 12.625) de \$7.477.000, que incluye \textbf{\$200.000 para el Puente Río Caldera} y \$90.000 para el tramo Castellanos--Yacones. Dos obras del departamento están en ejecución con certificados parciales. La \textbf{Ruta Nacional 55} pasa a cruzar Vaqueros y Calderilla hasta Tres Cruces. \\
1941 & \textbf{Año leído edición por edición} (ediciones 1878 a 1929). El expediente de pensión de la viuda de \textbf{Paulino Fernández}, «\textit{extinto agente de policía de la comisaría de la Caldera}», consigna que murió «\textit{en el punto denominado “la Calderilla” de resultas de un tiro de revólver}». \textbf{Primera muerte violenta que el archivo temprano registra en el departamento}, y llega por un trámite administrativo. \textbf{No consta el resultado judicial.} \\
1941 & Se declara «\textit{imprescindible}» adquirir «\textit{los terrenos necesarios}» para las mejoras del \textbf{camino Calderilla a río de las Pavas, tramo Desmonte--Campo Santo}. \textbf{Primera vez que una traza vial del departamento exige comprar el suelo}; no se publica de quién ni por cuánto. \\
1941 & Tras el remate desierto de mayo, un acto del mismo año \textbf{dispone inscribir la mina «Caldera nº 1, 2 y 3», expediente 44-L, «en calidad de vacante»}. En 1943 se le anula la patente. \textbf{Treinta y tres años desde el primer cateo de 1908 hasta la palabra «vacante».} \\
1941 & Se rematan las minas \textbf{«Caldera Nº 1, 2 y 3»} --- veintisiete hectáreas de plomo y plata en la finca San José, de Juan Larrán --- por \textbf{caducidad de la concesión ante la falta de pago del canon minero}. Base \$2.500. Nunca habían sido mensuradas. Decreto 4769 y edicto en el B.O.\ Nº 1.899 (30/05/1941); el edicto fija la subasta para el «sábado» 17 de julio, que fue jueves. \\
17 dic.\ 1941* & \textbf{Decreto 3012}: reelige a \textbf{Augusto Regis}\index{Regis, Augusto} Presidente de la Comisión Municipal «por un nuevo período legal». \\
29 dic.\ 1941* & \textbf{Decreto 3072-G}, expte.\ 1766-M/941: aprueba el balance del ejercicio 1940 de la Comisión Municipal. \textbf{Sin cifras}: es el primero de los cinco actos consecutivos sobre la hacienda del municipio que no publican un número. \\
1942 & \textbf{Año leído edición por edición.} Una inscripción registral transcribe el \textbf{acuerdo de la Municipalidad de La Caldera del 4 de junio de 1885} que concedió a \textbf{Eustaquio Murúa}\index{Murúa, Eustaquio} \textbf{treinta litros por segundo del río Wierna} para la finca «Sausal o Curuzú», en Vaqueros. Otra reconoce derechos sobre el río de las Nieves «\textit{otorgados por la Municipalidad de la Caldera}» y «\textit{desde fecha inmemorial}». \textbf{1942 es el año en que los derechos que el municipio concedió pasan a los registros provinciales.} \\
1942--1943 & \textbf{Augusto Regis}\index{Regis, Augusto} es nombrado \textbf{comisario de policía de La Caldera} mientras preside la Comisión Municipal ---reelecto por decreto los dos años--- y ocupa una banca de \textbf{diputado provincial}. Renuncia a la comisaría en 1943. \textbf{Undécima superposición de la serie, y la más completa.} \\
1942 & El P.E.\ \textbf{reconoce} el derecho de 1885 y ordena \textbf{inscribirlo en los registros provinciales}: un derecho otorgado por el municipio pasa a constar sólo en la Provincia. Se reconoce además \textbf{derecho exclusivo} sobre el arroyo del Abra. \\
1942 & Abre la estación de aforo \textbf{El Angosto}, sobre el río Mojotoro: \textbf{cuarenta y cuatro años de extensión}, de los cuales treinta y cinco quedan completos, y 828 km² de área de aporte. Cierra en 1986, lejos de los cincuenta años que el Decreto 1.989/02 exigirá. \\
5 may.\ 1942* & \textbf{Decreto 3829-G}: nombra a \textbf{Augusto Regis}\index{Regis, Augusto} \textbf{Comisario de Policía de La Caldera}. Presidente del municipio en 1933--1935 y de nuevo desde 1937, reelecto en diciembre de 1941, acumula desde aquí los dos cargos, hasta renunciar al policial en septiembre de 1943. \\
23 nov.\ 1942* & \textbf{Decreto 6749-H}: la Provincia paga \$200 a la Nación por \textbf{cinco mil alevinos de trucha arco iris sembrados en los ríos Castellanos y Wierna}. El decreto autoriza el gasto después del hecho y no dice en qué tramo ni en qué fecha se sembraron. \\
1943 & El \textbf{Gobierno de la Provincia} promueve \textbf{posesión treintañal} del terreno del pueblo «\textit{donde se encuentra emplazado el edificio ocupado por la comisaría de policía}», 40 por 200 metros, lindando con la calle principal y con Daniel Linares\index{Linares, Daniel}; se oficia a la Municipalidad «para que informe si afecta propiedad municipal». \textbf{Ocupación sin título desde 1913 o antes.} \\
1943 & El \textbf{Gobierno de la Provincia} deduce \textbf{juicio de posesión treintañal} sobre los ocho mil metros cuadrados donde funciona la Comisaría de Policía del pueblo: ocupaba sin título desde 1913 o antes, y \textbf{tuvo que ir a un juez}. \\
1943 & El juicio de posesión treintañal del Estado le cuesta \textbf{ochenta y dos pesos}: setenta al diario por el edicto y doce al Juez de Paz por diligenciar un oficio. \\
1943 & Se \textbf{anulan las patentes} de las minas ``Caldera Nº 1, 2 y 3'': el remate de 1941 \textbf{quedó desierto}. \\
1943 & \textbf{Augusto Regis\index{Regis, Augusto}} renuncia al cargo de \textbf{Comisario de Policía} de La Caldera. Es el último de los cargos que el archivo le registra desde 1908; el apéndice \ref{cap:personas} los enumera. \\
11 may.\ 1943* & \textbf{Decreto 5905-G}: reelige a \textbf{Augusto Regis\index{Regis, Augusto}} Presidente de la H.\ Comisión Municipal del Distrito de La Caldera «por un nuevo período legal» (arts.\ 178 y 182 de la Constitución). Cuatro meses después renuncia como Comisario de Policía del mismo departamento: \textbf{en 1943 los dos cargos recaen en la misma persona}. \\
14 jun.\ 1943* & \textbf{Decreto 6148-G}, expte.\ 190-M/943: aprueba la Ordenanza General de Impuestos \textit{y} la de Presupuesto de Gastos y Cálculo de Recursos de La Caldera para 1943. \textbf{No publica una sola cifra}; devuelve el expediente a la comuna. \\
4 ago.\ 1943* & \textbf{Decreto 208-G}: se reorganizan las Comisarías Inspectoras de Zonas. La Caldera queda en la \textbf{Primera Zona}, junto con La Capital, Campo Santo y Cerrillos: la única de las siete \textbf{sin comisaría inspectora propia ni asiento}, atendida directamente por Jefatura de Policía. \\
27 ago.\ 1943* & \textbf{Decreto 520-H, art.\ 22}: medio día de viático, \$6, al Jefe de la Sección Irrigación por la inspección realizada a \textbf{la finca Lesser}. El decreto no nombra departamento ni el objeto de la inspección. \\
30 sep.\ 1943* & \textbf{Decreto 755-G}, expte.\ 1549-M/943: aprueba el balance de la renta municipal de 1942, elevado conforme el art.\ 89 de la Ley 68. \textbf{Tampoco publica cifras.} \\
27 oct.\ 1943* & \textbf{Decreto 1066, art.\ 9º}: confirma a \textbf{Teófilo Reyes\index{Reyes, Teófilo}} como Juez de Paz Propietario del Distrito Municipal de La Caldera, da por terminadas las funciones de David Villa Gómez como suplente y nombra a Justo Robles. Los considerandos invocan «el movimiento revolucionario del 4 de junio». \\
30 oct.\ 1943* & \textbf{Decreto 1091-H} (B.O.\ Nº 2.028, 12 de noviembre, verificado sobre el facsímil): la Comisión Consultiva de Hidráulica contrata por licitación privada \textbf{defensas y arbolado sobre el río Mojotoro} por \textbf{\$24.100}, «\textit{por cuanto se acerca el período de lluvias}». El decreto no dice en qué tramo del río. \\
1944 & \textbf{Año leído edición por edición.} El \textbf{Revalúo General} se publica en el Boletín Nº 2166: es un padrón \textbf{provincial} de miles de partidas en el que «\textit{departamento de la Caldera}» abre una sección. Entre sus titulares figura \textbf{«Caldera, Dominga M.\ de»}. \\
1944 & Un decreto aprueba el presupuesto municipal «\textit{visto este expediente en el que el señor \textbf{Interventor de la comuna de la Caldera} eleva}» el proyecto de ordenanza. \textbf{Tercera intervención del cuerpo municipal} tras la de 1923 y la de 1930, y \textbf{ningún acto del año dice cuándo se dispuso ni quién es el Interventor}. El mismo decreto lo llama «comuna» y «Municipalidad». \\
1944 & Se licitan y ejecutan las \textbf{obras de refección de la comisaría de policía del pueblo}: presupuesto \$4.363,50, contratista Francisco Crescini, certificado final \$5.018,02. \textbf{Un año después del juicio de posesión treintañal por ese mismo terreno}: primero el título, después la obra. \\
1944 & El Jurado de Valuaciones publica el \textbf{Revalúo General}: cuarenta y siete partidas para todo el departamento, \$353.100 en total. \textbf{Francisco S. Urquiza\index{Urquiza, Francisco S.} es el mayor de los cuarenta y siete titulares publicados} con \$52.400. \\
19 abr.\ 1944* & \textbf{Decreto 2857-G}, expte.\ 6463-M/944: aprueba el Cálculo de Recursos y Presupuesto de Gastos de 1944, elevado ya no por un Presidente sino por \textbf{«el Interventor de la Comuna de La Caldera»}. Sin cifras. \textbf{El municipio fue intervenido entre junio de 1943 y abril de 1944}; el acto de intervención no aparece en el barrido. \\
4 may.\ 1944* & \textbf{Decreto 3016-H}, expte.\ 15.993/1944: aprueba la eliminación del \textbf{camino de Vaqueros a Los Yacones} del plan de obras de la Ley nacional 12.776 y reasigna sus \textbf{\$40.000} al camino de Salta a La Floresta, fuera del departamento. El fundamento consta: el gasto «no estaría justificado con el beneficio que la inversión reportaría». \\
20 sep.\ 1944* & \textbf{Decreto 4614-H}, expte.\ 17.870/1944: concede a los \textbf{Ferrocarriles del Estado} cinco litros por segundo del \textbf{río Mojotoro} ---el pedido hablaba de «un total diario de 100.000 litros», cifra que no concuerda con ese caudal--- para las locomotoras entre Güemes y Salta. Carácter \textbf{precario} «hasta tanto entre en vigencia el decreto ley de agua de la Provincia»; la Provincia se exime de responder por «la existencia en las diferentes épocas del año» del caudal. \textbf{La toma se remite a un plano adjunto que no se publica.} \\
1945 & \textbf{Año leído edición por edición.} \textbf{Tercer intento de defensa del río de la Caldera}: la Dirección General de Hidráulica llama a \textbf{licitación pública} con presupuesto de \textbf{\$2.515} ---aviso repetido en dieciséis ediciones---, el decreto 8549 del 3 de septiembre autoriza la ejecución, y en noviembre se resuelve hacerla \textbf{«por vía administrativa»}. \textbf{Ni certificado ni recepción se publican.} La refacción de la comisaría del año anterior había costado \$5.018. \\
1945 & Se traslada la escuela Nº 250 de \textbf{Wierna} «\textit{por encontrarse despoblado dicho lugar}» a San Francisco, partido de Yacones; y se aprueba una escuela de la Ley 4874 en \textbf{San Alejo}, «\textit{lugar donde según censo practicado existe una población de \textbf{51 niños} en edad escolar}». \textbf{Primer dato de población de un paraje que el archivo temprano publica.} \\
1945 & La licitación pública de las obras de defensa del río \textbf{queda desierta}: no se presentó ningún proponente. Se ejecuta \textbf{por administración} para llegar antes de las lluvias, por \$2.904,83 contra los \$2.515 licitados. \\
3 ene.\ 1945* & Edicto Nº 460, veintitrés ediciones: \textbf{posesión treintañal de la finca «Las Lagunas»}, departamento La Caldera, por Gregorio Maidana. Descripta por accidentes ---arroyo Los Comederos o El Barbecho, serranías de La Caldera, \textbf{río de Los Yacones}, arroyo Las Carretas---: \textbf{sin superficie, sin catastro, sin coordenadas}. \\
10 ene.\ 1945* & \textbf{Decreto 5738-G}: fija el \textbf{río Vaqueros} como límite norte de la jurisdicción policial de la Capital, y excluye del Destacamento Norte a la \textbf{Destilería Chachapoyas}, «que posee jurisdicción policial propia». El decreto da por existente esa jurisdicción y no cita la norma que la crea. \\
12 ene.\ 1945* & \textbf{Decreto 5786-G}, expte.\ 8266/944: aprueba el presupuesto municipal de 1945 «\textit{que corre a fojas 19 y 20}» del expediente, \textbf{sin transcribirlo}. Es el quinto acto consecutivo sobre la hacienda de La Caldera sin una sola cifra publicada, y el más rápido de toda la serie. \\
21 jun.\ 1945* & Remate judicial Nº 875: \textbf{cuarta parte indivisa de la finca «Chalchanio»}, 490 ha 9 a 36 cm², base \$671,01, en autos \textit{Strachan, Yáñez y Cía.\ c/ Martín Herrera}. \textbf{Mismo folio 342, asiento 438 del Libro B} que el edicto de 1937, con otro ejecutado y con la finca ya partida entre comuneros ---uno de ellos, Ernesto Echenique---. El aviso declara vigencia hasta el 27 de julio y deja de publicarse el 7. \\
3 jul.\ 1945* & Auto del juzgado; edicto Nº 1214, treinta ediciones desde el B.O.\ Nº 2.403: \textbf{Hermann Pfister\index{Pfister, Hermann}} promueve posesión treintañal sobre un terreno en \textbf{La Calderilla} de 108,28 m de frente por 2.395 m de fondo, con el \textbf{río de La Caldera} como lindero oeste. El perito de los deslindes de 1932, 1934 y 1937 aparece como parte. \\
15 dic.\ 1945* & \textbf{Decreto 9686-G}: convoca a elecciones nacionales y provinciales para el \textbf{24 de febrero de 1946}. La Caldera integra el grupo de dieciséis departamentos que eligen \textbf{un diputado y un senador} titulares, el escalón más bajo de los tres; La Capital elige siete diputados. \\
1946 & \textbf{Año leído edición por edición; es el último tramo con lectura sobre la imagen.} \textbf{Distribución de caudales del río Mojotoro}: considera que el Mojotoro forma su caudal con los ríos \textbf{Wierna, la Caldera y Castellanos}, divide el sistema en zonas ---la II abarca «\textit{desde la confluencia de los ríos la Caldera y Vaqueros}»---, fija un umbral de \textbf{2.500 litros por segundo} en la zona I y dispone que «\textit{las municipalidades de Campo Santo y de la Caldera prestarán a la Dirección General de Hidráulica la \textbf{colaboración necesaria}}». \textbf{Conceder (1885), reglamentar (1914), integrar una comisión (1931), colaborar (1946).} \\
1946 & Se publica el \textbf{balance de la comuna} del cuarto trimestre de 1945, \textbf{con cifras}: saldos de \$98,15 y \$98,02, firmado por \textbf{F.\ Mercado}. \textbf{Primer balance municipal con números desde 1908}, y lo firma un interventor. En abril Mercado renuncia; el 6 de agosto un decreto \textbf{da por terminada la intervención} y designa presidente a Enrique Alejandro Pfister. \textbf{La tercera intervención duró al menos tres años.} \\
6 jul.\ 1946* & \textbf{Decreto 659-G}: da por terminada la intervención a la Municipalidad de La Caldera y nombra presidente de la Comisión Municipal a Enrique Alejandro Pfister (B.O.\ Nº 2.608). \\
16 oct.\ 1946* & \textbf{Decreto 1966-H}, en acuerdo de ministros, a pedido de los regantes y de la Municipalidad de Campo Santo: mientras el Wierna lleve menos de 2.500 l/s, las fincas de la Zona I (río Wierna, La Caldera y afluentes) no pueden derivar más de \textbf{382 l/s}; usaban 1.123. Las municipalidades de Campo Santo y La Caldera deben colaborar con la Dirección General de Hidráulica (B.O.\ Nº 2.686). \\
18 dic.\ 1946* & \textbf{Decreto 2599-H}: reconoce los servicios de dos \textbf{jueces de distribución de las aguas} de los ríos Wierna y Mojotoro, pagados por la Provincia (B.O.\ Nº 2.736). \\
1947 & \textbf{Año barrido edición por edición al cierre} (ediciones 2869 a 3026), \textbf{con pre-informe automático y sin lectura sobre la imagen} (apéndice \ref{cap:fuentes}). La \textbf{Ley Nº 853}, publicada el 7 de agosto (verificada sobre el facsímil), fija la delimitación territorial de todos los distritos municipales de la provincia y dice que «\textit{El Municipio de La Caldera abarca todo el territorio del Departamento de su mismo nombre}»: \textbf{en 1947 hay un solo municipio y coincide con el departamento}, con la misma fórmula que la ley usa para Cachi, Cafayate y San Antonio de los Cobres. Cinco ediciones publican el edicto de reconocimiento de \textbf{50 litros por segundo del río de la Caldera} para la finca \textbf{«La Helvecia»}, 24 hectáreas. El \textbf{decreto 5276-E del 30 de julio} publica la tabla que fija los porcentajes de coparticipación municipal: por el art.\ 6º de la \textbf{Ley 834} se liquida «\textit{en proporción a los recursos efectivamente percibidos por cada una durante el año 1946}». \textbf{La Caldera recaudó \$3.068 en 1946 y le corresponde el 0,127 \%}; la capital, \$1.298.124 y el 54,116 \%. Un decreto de septiembre reparte la participación de la \textbf{Ley nacional 12.956} entre \textbf{42 municipalidades} por \textbf{\$209.775,06}: a La Caldera le tocan \textbf{\$266,41}, el \textbf{0,127 \%}, contra el 2,38 \% que daría un reparto parejo. La \textbf{Ley 834} de presupuesto (separata del 19 de mayo, verificada) fija el presupuesto provincial en \$13.560.293,76 y destina en su \textbf{Ítem 6} \textbf{\$100.000 al Mojotoro, \$50.000 al río La Caldera y \$50.000 al Wierna}, sobre un total de \$560.000 para defensas contra crecientes en toda la provincia: \textbf{el 35,7 \% va a los tres ríos del departamento}. Su \textbf{Inciso VI}, con fondos de la \textbf{Ley 770}, destina \$370.000 a estaciones sanitarias en seis localidades, entre ellas La Caldera. Y su \textbf{artículo 6º} fija el reparto de la coparticipación municipal. El \textbf{decreto 5466-E del 14 de agosto} aprueba el \textbf{plan hidráulico de obras 1948--1950} de la Administración General de Aguas por \textbf{\$35.100.000}, y \textbf{no lo publica}: el plan corre «\textit{de fs.\ 12 a 18}» del expediente 17847/1947. El decreto \textbf{5537-G del 20 de agosto} adjudica \textbf{a la Cárcel Penitenciaria} la mesa escritorio y las tres sillas del Registro Civil del pueblo, por \$132,40. El decreto \textbf{4945 del 1º de julio} destina \$600.000 a mercados en municipios, con un \textbf{Mercado Tipo I para La Caldera}; en diciembre se lo \textbf{desafecta} junto con los de La Candelaria, La Poma, San Carlos y San Antonio de los Cobres, para financiar el Mercado Frigorífico de \textbf{Embarcación}, «\textit{a mérito de considerarse ineficiosa la construcción \ldots{} expresión que se había manifestado por los propios habitantes}» (verificado sobre el facsímil). El \textbf{decreto 5916-E del 22 de septiembre} (verificado sobre el facsímil) aprueba el Acta 188 del Consejo de Vialidad \textbf{«a excepción de la resolución Nº 4967»}, y le exige a Vialidad «\textit{informar en mérito a qué consideraciones ha sido incluido, \textbf{fuera del plan}, el camino denominado «La Caldera a Santa Rufina», en la red de caminos provinciales}». El plan de obras de la Administración de Vialidad consigna «\textit{Puente S/Río La Caldera. Obras de arte Luz 150 mts.}» por \textbf{\$690.000}, sobre un plan total de \$12.867.275,27 (verificado sobre el facsímil). \\
1947 & Funcionan \textbf{cinco estaciones de aforo} en la cuenca, \textbf{las cinco de la Subsecretaría de Recursos Hídricos de la Nación}. La única estación provincial de la cuenca será la pluviométrica de Los Yacones, entre 1972 y 1990. Todas cerradas hoy; sus registros están \textbf{publicados} en compilaciones nacionales de 2004 y en el SNIH. \\
24 jul.\ 1947* & \textbf{Ley 2131 (original 853)}: fija los límites de los distritos municipales de la provincia. «El Municipio de la Caldera abarca todo el territorio del Departamento». Su art.\ 2 dispone que, hasta un censo general, los municipios mantienen sus categorías y organización (B.O.\ Nº 2.912, 07/08/1947). \\
1948 & \textbf{Año barrido edición por edición al cierre} (desde la edición 3060), con la misma reserva. \textbf{Decreto 8176 del 5 de febrero} (verificado en el sumario de la edición 3062, h.\ 3): fija ubicación y circuitos de los colegios electorales 1 a 6, y \textbf{La Caldera es el Colegio Nº 2}, con circuitos 5 y 6. El del circuito 5 trae la \textbf{decimotercera nómina de parajes}, la primera posterior a 1940. \textbf{Decreto Nº 9041-E del 5 de abril} (verificado sobre el facsímil): \textbf{Hermann Pfister} ---el perito que deslindó Loma Bola, Santa Mónica, Quitilipe y el Potrero Gallinato en 1916 y 1917--- pide \textbf{50 l/s} del río La Caldera para la finca \textbf{«La Helvecia»}, 286 ha 8.608 m² con 24 regables (expte.\ 4887/A/1948), y \textbf{se le reconocen 12,6 l/s}, con carácter temporal y permanente, por el art.\ 254 del Código de Aguas. La \textbf{Municipalidad de La Caldera}, consultada por la Administración General de Aguas por el inciso a) del art.\ 350, informa «\textit{que no tiene observación alguna que formular}». En los comicios de constituyentes del \textbf{5 de diciembre}, el Partido Peronista propone por el departamento a \textbf{Alberto Ovejero Paz y Alejandro Pfister} (Tribunal Electoral, aviso 4375). Por \textbf{Decreto 8278-G del 14 de febrero} (verificado sobre el facsímil) el gobernador nombra «\textit{en carácter de reelección}» a \textbf{Enrique Alejandro Pfister} presidente de la Comisión Municipal, por el art.\ 178 de la Constitución y el art.\ 35 de la \textbf{Ley Nº 68}; es uno de \textbf{seis decretos consecutivos del mismo día} sobre seis municipios. \textbf{El circuito Nº 6 del Decreto 8176 es Vaqueros}, y su nómina encabeza con la finca \textbf{Entre Ríos}. \textbf{Daniel Linares} ---el propietario de Getsemaní en el deslinde de 1916--- pide agua del río para «\textit{El Angosto}», 29 hectáreas (expte.\ 4498/47), y se le reconocen \textbf{15,22 l/s} con carácter \textit{temporal y eventual}. \textbf{La razón caudal/superficie es idéntica a la de Pfister: 0,525 l/s por hectárea}, de modo que el criterio, aunque no se enuncia, se deduce de los dos actos. La \textbf{Resolución de Minas 591} del 2 de abril (verificada sobre el facsímil) adjudica la mina de \textbf{plomo y plata} «Caldera Nº 1, 2 y 3», expte.\ \textbf{44-L de 1929} ---la que el padrón de 1933 inscribía a nombre de Juan Larrán---, \textbf{a Enrique A.\ Zuccarino y María Josefina Méndez Alcorta de Zuccarino}: diecinueve años entre la inscripción y la adjudicación, y un traspaso que el Boletín no publicó. Y el \textbf{Decreto 8706-E del 10 de marzo} designa receptor de rentas del departamento a \textbf{Gerardo Guerrero}, «\textit{atento a lo solicitado por el Banco Provincial de Salta}». \textbf{\$35.026,99} para los sanitarios de la escuela provincial del pueblo. \\
1949 & \textbf{Año barrido edición por edición al cierre} (281 ediciones, 4.869 hojas), con la misma reserva. \textbf{Decreto 13793-E del 31 de enero} (verificado): certificado Nº 1 por \textbf{\$27.024,10} de trabajos ejecutados en la \textbf{Estación Sanitaria de La Caldera}, adjudicada a \textbf{E.C.O.R.M.\ S.R.L.} por decreto 10668 del 29 de julio de 1948; la obra se entrega ese año. \textbf{En la misma hoja} se certifica el \textbf{Mercado Tipo II de Pichanal}, adjudicado por el \textbf{decreto 7321 del 13 de diciembre de 1947} ---el mismo que desafectó el mercado de La Caldera---. El \textbf{Ítem 6} del presupuesto reparte \$460.000 y el \textbf{río La Caldera se lleva \$80.000}, la partida más alta de las dieciocho; \textbf{Vaqueros aparece como río propio por primera vez}. La \textbf{Resolución de Minas 672 del 10 de febrero} vuelve a adjudicar la mina «Caldera 1, 2 y 3» a Enrique A.\ Zuccarino y María Josefina \textbf{Méndez} Alcorta de Zuccarino, \textbf{lugar denominado San José}, con el mismo texto que la 591 de 1948. \\
1950 & \textbf{Año barrido edición por edición al cierre} (277 ediciones, 4.156 hojas). Se cierra la \textbf{estación sanitaria}: se devuelven los certificados «\textit{con imputación a la cuenta especial "Depósitos en garantía"}», que es la recepción definitiva, cuatro años después de haber sido presupuestada. Tres ediciones publican el edicto de reconocimiento de agua sobre la propiedad \textbf{«Fracción Entre Ríos», ubicada en Vaqueros}: \textbf{la finca del municipio ya está dividida}. Y la comisión organizadora de los \textbf{festejos patrios} del pueblo pide un subsidio de \$500. \\
1951 & La \textbf{A.G.A.S.} proyecta defensas sobre el río La Caldera por \$27.781,95 y tres defensas de rama y piedra sobre la margen izquierda del Vaqueros. Ambas se imputan a la partida \textbf{``obras de defensa permanente''}. \\
1951 & \textbf{Siete actos sobre derechos de agua} del departamento en un solo año, casi todos \textbf{reconocimientos} de derechos preexistentes. El municipio no interviene en ninguno. \\
24 ago.\ 1951* & Sanción de la Ley Nº 1338, renumerada 2616 (promulgada el 27, B.O.\ Nº 4.025 del 3 de septiembre): autoriza a enajenar tierra fiscal por adjudicación directa, con preferencia para los actuales ocupantes. \\
1970* & Ley 4365: creación del municipio de Vaqueros. Primer intendente, Julio Catalán Arellano. Sede en terreno donado por la sucesión del Dr.\ Carlos Serrey\index{Serrey, Carlos}. \\
1972* & Inicio de construcción del embalse Campo Alegre. Adjudicataria: Rodio Argentina S.A. \\
8 jun.\ 1972* & Ley provincial Nº 4486: declara de utilidad pública y expropia 686,19 ha en el departamento La Caldera para el Embalse de Campo Alegre y el desvío de la RN 9. Afecta a Campo Alegre (Isaac Fernández), Los Porongos (Jorge Washington Álvarez), El Angosto o Angostura (Jaime Durán) y dos fracciones de la familia Mercado, que aportan el 53,2\% del total. Compensación fijada en valuación fiscal más 30\%, con toma de posesión previa a la discusión del precio. Sancionada y promulgada por el Ejecutivo de facto en un solo acto. \\
24 oct.\ 1973* & B.O.\ 9.371: Decreto 1.410/73, reglamentación de loteos urbanos y suburbanos. Obras antes de la venta; matrícula después de la verificación municipal. Regirá 46 años. \\
23 nov.\ 1973* & Ley Nº 4708: incluye en el Plan de Obras Públicas 1974 el cordón cuneta y la pavimentación de la calle principal de La Caldera y la ruta de acceso al Cristo Redentor. Promulgada por el gobernador Miguel Ragone. \\
1980 & \textbf{Decreto 12}: se declara \textbf{zona de emergencia agropecuaria} el área tabacalera de La Caldera y Chicoana por daños de granizo. \textbf{En 1979 había tabaco en el departamento y el Estado lo reconocía por decreto.} \\
1980 & \textbf{Permiso precario de uso} al Club de Regatas Campo Alegre sobre hectárea y media del \textbf{perilago} del dique: revocable y por cinco años. En 2015, el mismo tipo de cesión será de 41 ha por veinte años. \\
1980 & En veintitrés días la Provincia aprueba \textbf{tres encauzamientos} --- ríos Wierna, La Caldera (zona Campo Alegre) y Vaqueros --- por \$527.621.981 y \textbf{los tres por vía administrativa}, sin licitación. Nominalmente, el doble del presupuesto municipal sancionado en 1979. \\
1980 & El \textbf{Observatorio Pluviométrico Los Yacones} funciona con observador pago por la Provincia: \textbf{había capacidad de medir}, aunque la Provincia no aforaba el río: las estaciones de aforo de la cuenca eran nacionales, y la de El Angosto atravesaba un tramo sin registro. \\
1980 & \textbf{Vialidad Nacional} licita \textbf{defensas de hormigón} sobre la Ruta 9, tramo Río Caldera--Alto de la Sierra, por \$62.264.460 y cuatro meses de plazo. Y el Plan de Trabajos Públicos asigna \$300.000.000 al \textbf{proyecto de la potabilizadora del Dique Ing.\ José Alfonso Peralta}. \\
30 jun.\ 1983* & Ordenanza 047 de La Caldera, firmada por \textbf{Arturo R. Fernández, Presidente de la Comisión Municipal}, y aprobada por decreto del Ministerio de Gobierno. El Decreto 877/82 había concentrado en el ejecutivo las funciones del Concejo. \\
30 jul.\ 1983* & Ley 6131: convoca a los comicios del 30/10/1983. La Caldera elige un diputado y un senador. Los municipios eligen \textbf{concejales y miembros de comisiones municipales}: el ejecutivo municipal no está en la boleta. \\
20 oct.\ 1983* & Ley Nº 6181: expropia una hectárea del catastro 132 (Hilario Arias) para la Escuela Nº 660, Finca Mojotoro, por 12,49 pesos argentinos, valor fiscal más 30\%. Dictada bajo facultades legislativas de la Junta Militar. \\
9 dic.\ 1983* & Víspera del traspaso democrático: el P.E.\ provincial acepta las renuncias de los presidentes de las Comisiones Municipales, entre ellas las de \textbf{La Caldera y Vaqueros}. \\
3 sep.\ 1985* & Ley 6334: expropia 2 ha 1.736,80 m² de la Finca Vaqueros (catastro 1229, dominio de Manuel Serrey y otros) para viviendas; donación al IPDUV, con reversión si no construye en dos años (B.O.\ 12.311). \\
4 mar.\ 1986* & Ley Nº 6354: expropia 8.158,71 m² de Elena Serrey\index{Serrey, Elena} de González Bonorino en la localidad de La Caldera (matrícula 1.303) para construcción de viviendas familiares, con donación al IPDUV y cláusula de reversión a los dos años. \\
14 oct.\ 1986* & Decreto de necesidad y urgencia 2876/86: expropia 3.710,08 m² de la Finca Vaqueros (catastro 1280), de la sucesión testamentaria de Carlos Serrey, para la planta potabilizadora de Vaqueros. Vencido el plazo legislativo, el Decreto 1500/87 lo tiene por Ley 6464 (B.O.\ 12.757, 03/08/1987). \\
Años 90 & Resolución del juicio sucesorio de los descendientes de Serrey\index{Serrey, Carlos}: regulariza dominios y habilita la subdivisión y venta de parcelas en Vaqueros. \\
sep.\ 1996* & Ley 6895: expropia terrenos del departamento Capital sobre la margen derecha del río Vaqueros para venderlos a sus ocupantes y encarga defensas aguas abajo del puente de la Ruta 9; el Decreto 1962/96 veta la expropiación principal (matrícula 91.676) y la parcelación (B.O.\ 15.006). \\
nov.\ 2001* & Se aprueba y adjudica el \textbf{encauzamiento del arroyo Urquiza\index{Urquiza, Francisco S.}}, en Vaqueros, por \$7.356,80 y diez días de plazo, por vía abreviada. \\
2002 & \textbf{Decreto 1.989/02}: se publica como separata el procedimiento de determinación de línea de ribera. Doce artículos. Los \textbf{aforos con series de cincuenta años} encabezan la prelación de fuentes; la resolución se registra \textbf{previa publicación por dos días}; el trámite puede iniciarse \textbf{de oficio}. \\
2002 & Se licita la concesión del \textbf{corredor metropolitano de transporte} entre Capital, La Caldera y Cerrillos: pliego \$6.000. \\
23 ago.\ 2007* & Ley 7460: donación a la Municipalidad de La Caldera del catastro 802, manzana 50 del barrio Jardín, para uso público y recreativo (B.O.\ 17.706). \\
21 oct.\ 2008* & Decreto Nº 4.414/08, emergencia habitacional en todo el territorio provincial y Programa «Comunidad Organizada». Ordena reglamentar la transferencia de tierras fiscales; el reglamento llega en 2016. Excluye a quienes ocupen inmuebles ilegalmente. \\
Desde 2010 & Reclamos vecinales en La Caldera por protección de ríos y entorno ante la Secretaría de Ambiente y el municipio. \\
2010 & Comienzan proyectos de barrios privados en la zona alta de Vaqueros; se menciona el loteo Las Vertientes, con suministro propio. \\
2010* & Censo Nacional: 7.763 habitantes; 2.088 hogares para 2.578 viviendas particulares, con 490 sin residentes, el 19,01\% del parque contra el 4,87\% provincial. El departamento registra la tasa de ocupación más alta (64,37\%) y la de inactividad más baja (32,39\%) de los veintitrés departamentos de Salta. \\
2010 & \textbf{Noviembre}: el \textbf{``Plan de Mínima de Defensas y Encauzamientos de Ríos''} interviene el río La Caldera y los arroyos Guaranguay, \textbf{Los Duraznos} y Los Manzanos, con 250 horas de máquina. \\
13 oct.\ 2010* & Se constituye \textbf{Jardín Celestial S.R.L.} (B.O.\ Nº 18.481, 01/12/2010), con objeto de cementerio parque en jurisdicción del departamento La Caldera. Socios: Héctor D'Andrea\index{D'Andrea, Héctor} 55\%, Carlos Miguel Calabró\index{Calabró, Carlos Miguel} 30\% y Renzo Daniel Palacios 15\%. \\
10 nov.\ 2011* & Ley 7699: el Poder Ejecutivo dona a la Municipalidad de La Caldera el inmueble matrícula 1846, «con el cargo de ser destinado exclusivamente al desarrollo integral de la zona» (B.O.\ 18.728). \\
20 nov.\ 2011* & Se releva ---en una modalidad que el propio expediente de 2023 describe de dos maneras: «centralizado», por el equipo técnico del INAI, en un informe, y «desde el área de descentralizada del INAI» en otro--- la ocupación tradicional de la \textbf{Comunidad Originaria Kolla Kondorwaira}: la \textbf{zona alta de la cuenca del río La Caldera} es territorio de ocupación actual, tradicional y pública. Aprobado por Res.\ INAI 617/2013. \\
2012 & Plan de Ejecución Metropolitano: La Caldera con 37,3\% de crecimiento intercensal, el mayor del área. Advierte sobre la conversión en base de asentamiento de sectores cerrados. \\
2012 & El \textbf{Decreto 2478/12} aprueba el \textbf{PDES 2030} con carácter definitorio. Declara que el crecimiento del área metropolitana fluyó \textbf{``espontáneamente, sin planificación''} y enumera como amenaza la \textbf{falta de tierra pública apta para urbanizar}. \\
2012* & La Secretaría de Recursos Hídricos diseña el \textbf{Plan de Monitoreo de Rutina} de la calidad del agua de la subcuenca Mojotoro. \\
2013* & Se inicia ante la Secretaría de Recursos Hídricos el expediente Nº 34-167.127/13, de determinación de línea de ribera sobre el inmueble matrícula Nº 4.366, Finca La Caldera o Getsemaní. Su comisión técnica (Res.\ 395/13) nombra al río La Caldera y al arroyo El Manzano; del mismo expediente saldrá en 2014 la Res.\ 383/14 sobre los \textbf{arroyos El Durazno y Guaranguay}. \\
2013--15* & El CFI encomienda el PDUA de La Caldera; según la presentación oficial, el trabajo comenzó en 2013. El informe identifica loteos privados en ejecución sin planificación y deja pendiente el estudio de cota máxima edificable. Su introducción y cierre nombran, por error de plantilla, al municipio de Rivadavia Banda Sur. \\
30 may.\ 2013* & DNU Nº 1595/13, luego Ley Nº 7781: declara la Emergencia Hídrica por escasez de agua en todo el territorio provincial por un año, atribuyéndola fundamentalmente al cambio climático, y ordena readecuar y ordenar todas las cuencas. Se origina en la Resolución 142/13 de la Secretaría de Recursos Hídricos. \\
10 oct.\ 2013* & Res. 292/13: comisión técnica para el arroyo Quintín, catastro 3756. Con la 348/13, la 395/13, las 6, 7 y 8 de 2014 y la 383/14 son siete de los actos que la serie registra sobre la cuenca en esos catorce meses, a los que se suma la 023/14; la serie 2013--2015 reúne veinticinco actos de comisión o determinación, más la 023/14, que designa y determina a la vez. \\
19 nov.\ 2013* & Res. 348/13, línea de ribera del río La Caldera, matrícula 1590, La Calderilla. Su edicto (B.O. 19.213, 23 y 26/12/2013) \textbf{publica las coordenadas Posgar}, determina la \textbf{zona inundable} y transcribe el art. 172 del Código de Aguas. \\
26 dic.\ 2013* & Res. 395/13: aprueba la comisión técnica para el río La Caldera y el arroyo del Manzano, matrícula 4366 ---el escaneo dice «1366»---, Finca La Caldera o Getsemaní (Expte. 34-167.127/13). \\
13 ene.\ 2014* & Res. 6/14, 7/14 y 8/14: dos determinaciones de línea de ribera sobre el río La Caldera (matrículas 3593 y 3584), cuyos avisos de enero publican la tabla de la línea del pueblo que aprobará la 023/14, y una comisión técnica para la quebrada de la matrícula 4191, Fracción Finca Vaqueros. \\
30 ene.\ 2014* & \textbf{Res. 59D/14}: aprueba el plano del Club de Campo ``Las Vertientes'', matrícula 3989 de Vaqueros, 143 ha. Condiciona las matrículas a reglamento de condominio con \textbf{servicio de agua privado}, \textbf{resolución de línea de ribera} y libre deuda. \\
feb.--mar.\ 2014* & Edictos de renovación de las canteras de áridos ``Macarena'' (Transporte Hnos. Mendoza S.R.L., río Wierna, 10 ha) y ``El Vaquero'' (INCOVI S.R.L., Vaqueros, 9 ha). En ambos casos, terrenos de propiedad fiscal. \\
11 feb.\ 2014* & Decreto 397: contrata al coordinador del \textbf{Plan Urbano Ambiental de Vaqueros}, en el marco del Programa de Mejora de la Gestión Municipal, \textbf{préstamo BID 1855/OC-AR}, excluido del régimen de la Ley 6.838. \\
19 feb.\ 2014* & Res. 023/14 (del 6 de febrero), a pedido del intendente Mendaña: designa al geólogo \textbf{Jorge G. Torres} y \textbf{determina la línea de ribera del río La Caldera a lo largo del pueblo}, con 134 coordenadas publicadas en el anexo vinculado al aviso; la sujeta a la construcción de \textbf{defensas con gaviones}, que tramitan por el expte.\ 34-253.934/13 (expte.\ de la resolución: 0090034-96837/2013). \\
11 mar.\ 2014* & Res. SOP 114/14: adjudica a Dal Borgo Construcciones la obra de \textbf{protección y defensas de las cámaras de captación del Acueducto Norte del río La Caldera}, \$1.824.588. Entre los oferentes admisibles, INCOVI S.R.L. \\
abr.\ 2014* & Renovación de la cantera de áridos ``El Triunfo II'' (Cooperativa La Mesa Redonda Ltda.), sobre el río La Caldera, departamento Capital. Tercera concesión de áridos de la cuenca en un trimestre. \\
7 abr.\ 2014* & Res. 80/14: aprueba la línea de ribera del arroyo Quintín o Pacará sobre la \textbf{matrícula 3989, fracción 78} --- el inmueble del Club de Campo Las Vertientes ---, sesenta y siete días después de que la Res. 59D/14 la exigiera. El edicto publica las coordenadas. \\
21 abr.\ 2014* & Res. SOP 220: aprueba el proyecto ejecutivo de la obra ``Acueducto, Mejoras en Planta Potabilizadora, Canal de Aducción y Redes Maestras'' para la localidad de Vaqueros, \$1.500.000, 180 días. \\
29 abr.\ 2014* & Res. 96/14 (comisión, río Vaqueros) y Res. 101/14 (determinación, arroyo Los Nogales Norte o El Cajón, matrícula 4191, \textbf{con coordenadas publicadas}). \\
26 may.\ 2014* & Acta de socios de Jardín Celestial S.R.L.: se amplía el objeto a la comercialización y adquisición de inmuebles y se eleva el capital a \$545.000 (B.O.\ Nº 19.394, 29/09/2014). \\
27 may.\ 2014* & Res.\ 142/14: determinación de línea de ribera de los ríos Mojotoro y La Caldera, matrícula 3616, \textbf{sin coordenadas publicadas}. Su aviso sale recién el 13 y el 14 de agosto (B.O.\ Nº 19.363 y 19.364). El 25 de agosto la Res.\ 236/14 la rectifica, también sin publicarlas (aviso del 3 de septiembre, B.O.\ Nº 19.377). \\
3 y 11 jul.\ 2014* & Res. 182/14 (río La Caldera, matrícula 2882, La Calderilla) y \textbf{Res. 185/14} (río La Caldera y arroyo El Manzano, \textbf{matrícula 4366}, Finca La Caldera o Getsemaní). Ambas con coordenadas publicadas. \\
29 jul.\ 2014* & Res. 209/14: determinación de la línea de ribera del río Vaqueros, matrículas 1905 y 3740, con coordenadas publicadas. \\
10 oct.\ 2014* & Audiencia pública ambiental por la \textbf{Urbanización Cerro de Buena Vista} (matrículas 1713 y 1905), en el Centro Deportivo Municipal de Vaqueros, con el expediente disponible en el municipio. \\
14 oct.\ 2014* & \textbf{Decreto 3069/14}: ratifica el contrato de promoción turística con La Caldera S.A.\ para ampliar su hostería. El contrato, que fija las exenciones, no se publica. \\
18 dic.\ 2014* & \textbf{Resolución Nº 383} de la Secretaría de Recursos Hídricos: aprueba la determinación de la línea de ribera realizada por la Comisión Técnica en \textbf{ambas márgenes de los arroyos El Durazno y Guaranguay}, inmueble catastro Nº 4.366, departamento La Caldera. Las coordenadas Posgar no se publican: quedan a disposición en la sede del organismo. \\
15 ene.\ 2015* & Se publica la Ordenanza Nº 000437 de Vaqueros (Reglamento de Fraccionamiento de Tierras), promulgada el 29 de diciembre de 2014 por el Decreto 78/14. El anexo no se publica en el Boletín; está en el digesto municipal. \\
21 y 22 ene.\ 2015* & Se publica el edicto de la Res.\ 383/14 en los Boletines Oficiales Nº 19.466 y 19.467, pág.\ 407, orden 100045883. Apelable ante el Tribunal de Aguas en diez días hábiles. \\
9 feb.\ 2015* & \textbf{Decreto 463/15}: segundo contrato de promoción turística, con Proyecto 1 S.A., para ampliar el Hotel La Caldera. Se publica el 19 de febrero sin el contrato, «para su consulta en oficinas». \\
8 jun.\ 2015* & \textbf{Decreto 1968/15}: la Provincia declara a su favor la prescripción adquisitiva del terreno de la Escuela Nº 4118 (Finca Potrero de Gallinato, matrícula 3162), con el plano 000962. \\
jul.\ 2015* & Se publica la gestión de agua subterránea por dren horizontal para \textbf{400 lotes del loteo El Durazno}, con carácter permanente, a nombre de una co-titular del catastro 4366. \\
10 y 12 ago.\ 2015* & Res.\ 579 D/15 (\textbf{Barrio La Misión}, libre deuda; el plano descuenta la ribera) y Res.\ 585 D/15 (\textbf{Chacras del Gallinato}, 1.569 ha, seis condiciones). \\
24 ago.\ 2015* & \textbf{Decreto 2949/15}: comodato gratuito por veinte años a Proyecto 1 S.A.\ de 41 ha fiscales para una cancha de golf. \\
oct.\ 2015* & El Chalchal S.R.L.\ gestiona agua subterránea de pozo perforado, con carácter permanente, para 580 habitantes del loteo La Misión Etapa III (catastro 3616). \\
30 oct.\ 2015* & Ordenanza 469/15 de Vaqueros: fija por coordenadas los ejidos urbanos de Vaqueros y de Lesser, «jurisdicción del municipio de Vaqueros». Decreto 112/15. \\
12 nov.\ 2015* & Ordenanza 474/15 de Vaqueros: Plan de Desarrollo Urbano Ambiental (PIDUA), trece días antes que el de La Caldera. Decreto 119/15. \\
25 nov.\ 2015* & Ordenanza Nº 643: se aprueba el PDUA de La Caldera. Mantiene vigentes las ordenanzas 040/82 y 243/97. \\
21 dic.\ 2015* & Decreto Nº 120/15: designa a Fernando Gabriel Martinis Mercado Secretario de Tierra y Bienes. \\
2016 & Vecinos autoconvocados advierten que el Acueducto Norte implicará el desmonte del cordón ribereño. \\
2016* & Sanción del Código de Ordenamiento Territorial de La Caldera, Ordenanza 06/2016. \\
7 jul.\ 2016* & Ordenanza 495/16 de Vaqueros: acepta la donación de calles (2 ha 4.333,21 m²), ochavas y 7.102,65 m² de espacios verdes del «Loteo Cerros de Buena Vista», matrículas 1713 y 1905. Decreto 147/16. \\
09 sep.\ 2016* & Se publican en el B.O.\ Nº 19.861 la Resolución Nº 82/16, que aprueba el Reglamento de Asignación de Lotes de Interés Social ---su anexo, sólo como imagen escaneada---, y la Nº 86/16, que modifica su art.\ 7º. \\
23 sep.\ 2016* & Acta Acuerdo del \textbf{Plan de Mínimas y Encauzamiento 2016--2017} (Decreto 1.889/16). La Provincia y \textbf{veintiún municipios}, entre ellos La Caldera y Vaqueros, con fondos de la Ley 7.931 de crédito público. \textbf{Los municipios ejecutan y cobran contra certificados.} \\
20 oct.\ 2016* & Ordenanzas 502/16 y 503/16 de Vaqueros: pasarela en el acceso a Villa Sara y ciclovía en Lesser. El intendente las promulga condicionando su ejecución a fondos provinciales o nacionales, por «la escasa coparticipación» del municipio. \\
7 nov.\ 2016* & Resolución Nº 360 de Recursos Hídricos: línea de ribera de ambas márgenes del \textbf{río Wierna}, siete matrículas del departamento, \textbf{sin coordenadas publicadas} (edicto en el B.O.\ Nº 19.951, 26 y 27/01/2017). Comisión: Torres y Cerezo. \\
17 nov.\ 2016* & Resoluciones Nº 7/16, 8/16 y 9/16 de la Secretaría de Tierra y Bienes: requisitos, baremo de puntaje y procedimientos de asignación (B.O.\ Nº 19.906). Los anexos de las tres se publican sólo como imágenes escaneadas; una fe de erratas sale en el B.O.\ Nº 19.911. \\
25 nov.\ 2016 & El Gobierno de Salta entrega la \textbf{carpeta técnica de relevamiento territorial} a la \textbf{Asociación Comunitaria Kondorwaira}, cinco años después del relevamiento de campo. \\
26 dic.\ 2016* & Ordenanza 507/16 de Vaqueros: aprueba el Ordenamiento Normativo Urbanístico ---Código de Planeamiento Urbano Ambiental, Código de Edificación y Lineamientos del Área Turismo---. Promulgada el mismo día por el Decreto 165/16. \\
2017* & Dos \textbf{contrataciones directas} de Vialidad con la Municipalidad de La Caldera por \$1.099.956: mantenimiento vial en cinco rutas provinciales. \textbf{La cotización es idéntica al presupuesto oficial} en los dos casos. \\
sep.\ 2017* & Constitución del fideicomiso Lares de la Inmaculada: 23 ha sobre la matrícula 3012, desarrollo a cargo de MDay. \\
31 oct.\ 2017 & La comunidad kolla \textbf{Kondorwaira} --- \textbf{27 familias} del paraje Potrero de Castilla, cacique \textbf{Teodora Sarapura} --- recibe un segundo aporte de \$1.028.000 del PRODERI para ganadería menor y producción de quesos. \\
2018* & Resolución Nº 21/2018 de Jefatura de Gabinete: crea una Comisión para tratar la problemática de loteos, urbanizaciones, clubes de campo y asentamientos de la provincia. \\
28 sep.\ 2018* & Ordenanza 535/18 de Vaqueros: manual de procedimiento para nuevas urbanizaciones en seis pasos, con licencia de comercialización, garantía y audiencia pública en instalaciones municipales; modifica parcialmente la 437. Decretos 206/18 y 207/18. \\
8 nov.\ 2018* & Acta de Jardín Celestial S.R.L.: Héctor D'Andrea\index{D'Andrea, Héctor} queda como Gerente Titular y \textbf{Héctor Miguel Calabró\index{Calabró, Héctor Miguel}} como Gerente Suplente (publicado en el B.O.\ Nº 20.675 el 30/01/2020). \\
21 dic.\ 2018* & Aluvión. Afecta El Durazno, El Nogalar y Potreros de La Caldera. \\
2019 & Los vecinos cuestionan la falta de monitoreo sobre la traza del acueducto, sobre franja ribereña de alta fragilidad. \\
2019 & Relevamiento universitario (Polliotto, Lema, Alonso Maurizzio y Reyes, UCASAL): 5 urbanizaciones privadas y 11 clubes de campo aprobados por la DGI, frente a 54 asentamientos. \\
3 ene.\ 2019* & Segundo aluvión. \\
feb.\ 2019 & La \textbf{Comunidad Kolla Kondorwaira} publica \textit{Nuestro tiempo redondo}, editado por el municipio: define su territorio como \textbf{la alta cuenca del Mojotoro y del río Las Nieves}, e incorpora a su dinámica \textbf{Vaqueros, La Caldera y el norte de la ciudad de Salta}. \\
feb.\ 2019* & \textbf{Se presenta ante el Juzgado de Minas el amparo ambiental colectivo \textit{Lazarte Vigabriel}} contra la Municipalidad de La Caldera y el Ministerio de la Producción, por los aluviones de diciembre y enero en la microcuenca del río La Caldera, dentro de la cuenca del Mojotoro (CUIJ 17-00023860-0). \\
4 abr.\ 2019* & Audiencia pública por Lares de la Inmaculada en el SUM del Colegio San Cayetano (convocatoria en el B.O.\ del 22, 25 y 26 de marzo; expte.\ 0090227-118038/2018-0). Impugnación previa rechazada. \\
nov.\ 2019* & La Sala III de la Cámara de Apelaciones en lo Civil y Comercial rechaza el amparo de los vecinos de Vaqueros. \\
25 nov.\ 2019* & Decreto Nº 1682/19: aprueba la reglamentación del Título V de la Ley 2.308 y deroga el Decreto 1.410/1973, vigente durante 46 años. Sus considerandos declaran que la seguridad jurídica de los adquirentes de lotes requiere mayor injerencia estatal y que debe garantizarse el efectivo cumplimiento de las obras de infraestructura a cargo privado. Crea la Unidad Coordinadora de Loteos, Loteadores y Desarrolladores Inmobiliarios. Su anexo, publicado con el decreto, exige certificados de no inundabilidad y de aptitud ambiental, habilita la preventa con caución y la matrícula anticipada, y reemplaza la verificación municipal por un visado con plazo. \\
24 jun.\ 2020* & Decreto Nº 399/20: crea el Plan ``Mi Lote''. Meta: 10.000 lotes, 40\% metropolitano. \\
jul.\ 2020 & Empleados municipales cortan la RN 9. El intendente declara que la coparticipación no alcanza para pagar sueldos. \\
oct.\ 2020* & Notas 101 a 105/20 de la Municipalidad: intima a los loteos a presentar planos, certificados de no inundabilidad y de aptitud ambiental, que no había recibido. \\
oct.\ 2020* & Resolución Nº 466 D/20: encomienda las competencias de esa Unidad Coordinadora a la \textbf{Secretaría de Tierras y Bienes}, el mismo organismo que administra la tierra fiscal. \\
26 nov.\ 2020* & Ordenanza 567/20 de Vaqueros: excepción al Código de Planeamiento a favor de un particular, por incumplir el retiro de seis metros de la zona R1 (art.\ 101, que la ordenanza atribuye al Código de Planeamiento y que es del de Edificación). \\
dic.\ 2020* & Clausura preventiva del loteo La Falda (matrícula 154, km 16--17,5 RN 9) tras denuncia del municipio. El personal técnico determina que el sector no es urbanizable según el COT; se indica la vía del Proyecto Especial. \\
fines 2021* & Acuerdo entre la Secretaría de Minería y Aguas del Norte (expte.\ 99078-2021): nueva zona de exclusión para canteras junto a las captaciones. \\
29 ene.\ 2021* & Decisión Administrativa Nº 50/21: designa a Mariano Puertas Jefe de Programa de Planificación Urbana. \\
jun.\ 2021 & Se concluye la primera etapa del acueducto hasta la cisterna El Huaico, con financiamiento CAF. La inauguración oficial es de octubre de 2021: planta alimentada por el embalse, 4.000 m³/h y 22 km de acueducto. Según Aguas del Norte, en diciembre de 2022 el acueducto que toma del dique todavía no operaba en forma regular. \\
jul.\ 2021 & Manifestación en la zona alta de Vaqueros por acceso igualitario al agua potable, tras quince años de reclamos. \\
15 ago.\ 2021* & Elección senatorial: Abilés no es reelecta; la banca pasa a Calabró. \\
27 sep.\ 2021* & Decreto Nº 844/21: traslada la Secretaría de Tierras y Bienes del Estado al Ministerio de Infraestructura; la DGI pasa a depender directamente del Ministro. \\
18 nov.\ 2021* & Sesión de despedida de la senadora Abilés tras ocho años de mandato. \\
26 nov.\ 2021* & Resolución Delegada Nº 646/21: aprueba el plano de mensura, desmembramiento y loteo Nº 1.164, proyecto ``El Durazno'', matrícula 6.558. \\
dic.\ 2021* & Sentencia en \textit{Lazarte Vigabriel} (Expte.\ EMI 23860/2020): declara la obligación estatal de protección y recomposición de la cuenca y ordena la confección del Plan de Manejo Integral de la microcuenca, en dos etapas y con base en el principio de progresividad. Funda en el art. 41 CN, la Constitución provincial y las leyes 7070, 7017 y 7543; introduce la cuestión climática con cita de la OC-23/2017 de la Corte IDH y la Observación General Nº 15 del Comité DESC. Criterio de infraestructura verde sobre gris. \\
2022* & Censo Nacional: el departamento La Caldera pasa de 7.763 a 12.299 habitantes y de 2.578 a 4.991 viviendas particulares. Cuarto departamento de mayor crecimiento demográfico del país, detrás de Añelo (Neuquén, 68,4\,\%), San Vicente (BA, 65,1\,\%) y General Rodríguez (BA, 63,7\,\%), según los resultados definitivos, cotejados en veintitrés de las veinticuatro jurisdicciones. El departamento registra 4.003 hogares para 4.991 viviendas particulares: 988 sin residentes, el 19,80\% del parque, contra el 8,55\% provincial; una de sus fracciones censales alcanza el 43,17\%. Provincia: +19,2\% de población en viviendas particulares y +45,3\% de viviendas. La Caldera y Vaqueros figuran como municipios de tercera categoría. \\
2022 & \textbf{Marzo}: se licita el \textbf{puente sobre el río Vaqueros y la Circunvalación Este}, cuatro kilómetros y cuatro carriles, por más de \$1.341 millones, \textbf{financiada por Vialidad Nacional}. Tres oferentes, entre ellos \textbf{DESA --- Defensa y Encauzamiento S.A.} El pedido de un segundo puente databa de 1935. \\
7 ene.\ 2022* & Minería inspecciona Tierra Gaucha: opera Supercemento con guías de tránsito de Geomix. En marzo paraliza extractores sin permiso en la zona de exclusión de Aguas del Norte (Resolución SMyE 20/2022, 591 ha). \\
16 mar.\ 2022* & Se inicia el incidente INC 23860/1 (``piezas pertenecientes'') en el amparo. Jardín Celestial S.R.L.\ figura como tercero en la carátula, con apoderado José Luis Criado. \\
3 may.\ 2022* & Decisión Administrativa Nº 366/22: designa a José Javier Haro Director General de Tierras Fiscales y Bienes del Estado. \\
24 may.\ 2022* & Resolución Delegada Nº 369 D: acepta la donación de 8.126,79 m² (3,05\% del total) ofrecida por Héctor D'Andrea\index{D'Andrea, Héctor} y Jardín Celestial S.R.L. \\
ago.\ 2022* & Plan de Monitoreo de Aguas Superficiales 2022 de la subcuenca Mojotoro: siete puntos en la microcuenca del río La Caldera. \\
sep.\ 2022 & Publicación de la bióloga y vecina sobre el descontrol del negocio inmobiliario; levantada por prensa nacional. \\
19 sep.\ 2022* & Registro municipal de loteos sobre la microcuenca: El Nogalar, La Milagrosa, dos etapas de El Durazno y el club de campo Potreros de La Caldera figuran como \textbf{irregulares}. \\
22 sep.\ 2022* & Providencia del Juzgado de Minas: se tiene presente lo manifestado por el actor sobre la extemporaneidad e inadecuación del informe de la Fiscalía de Estado y sobre la omisión de la Municipalidad de \textbf{Vaqueros}. \\
12 oct.\ 2022* & Inspección provincial y municipal en la margen derecha: construcciones, defensas precarias y tamices de áridos dentro de la línea de ribera; una retropala acordona en el cauce material del desmonte de El Durazno. El mismo día, Recursos Hídricos pide fondos al CFI para estudiar la cuenca. \\
19 nov.\ 2022* & Denuncia penal por un incendio en el catastro 1555, inmueble con sumario por desmonte ilegal. \\
2 dic.\ 2022* & Plano de Aguas del Norte «Relevamiento Acueducto Norte --- Líneas de ribera, cámaras acueducto y canteras». El informe que lo acompaña declara 10.822 usuarios del Acueducto Norte y 11.587 del Wierna, todos en la capital. \\
2023 & \textbf{Agosto}: la \textbf{Provincia demandada tiene que preguntarse a sí misma} si hay comunidades en la cuenca. Fiscalía de Estado pide los antecedentes territoriales para convocar a las comunidades a los talleres del Plan. Se ratifican \textbf{dos comunidades} en el departamento. \textbf{Septiembre}: constancias de consulta libre, previa e informada. \\
20 ene.\ 2023* & Nota interna de la Secretaría de Recursos Hídricos, incorporada al expediente de Fiscalía de Estado: el Plan de Manejo debe ajustarse «a la escala predial»; la calidad del agua para consumo es «responsabilidad de los usuarios». \\
30 mar.\ 2023* & El senador Calabró presenta el Expte. 90-31.780/2023, instituyendo la Reserva Natural Corredor Biológico Abra de Santa Laura sobre el talweg del río Porongal. Girado a comisión. \\
1 ago.\ 2023* & Minería paraliza a la Municipalidad de La Caldera por extraer áridos en el área de exclusión del acueducto. \\
7 ago.\ 2023* & La Municipalidad remite a Fiscalía de Estado la lista de loteadores, productores y funcionarios a convocar al taller. \\
16 ago.\ 2023* & Recursos Hídricos desaconseja prolongar el terraplén de la cantera La Serena junto al camping y cortar el monte bajo que lo protege. \\
30 ago.\ 2023* & Fiscalía de Estado pide a Asuntos Indígenas la consulta previa a las comunidades Kolla Kondorwaira y Kolla Urkuwasi. \\
31 ago.\ 2023* & Diseño del taller participativo del Plan de Manejo, previsto para el 29 de septiembre en el Centro Gaucho El Palenque; convoca a siete loteadores, tres productores de áridos y dos comunidades kolla. \\
12 sep.\ 2023* & Consulta libre, previa e informada en la Municipalidad de La Caldera, con seis organismos provinciales, el intendente y dos representantes de cada comunidad kolla, \textbf{incluida la Secretaría de Minería} (el Plan de Manejo la fecha el 23). Las comunidades plantean dudas sobre obras hechas sin autorización ni consulta. \\
19 sep.\ 2023* & Asuntos Indígenas tiene la consulta por concluida. \\
20 sep.\ 2023* & Dictamen 129/2023 de Recursos Hídricos: 275 usuarios de agua para riego en el departamento, criterio para convocar al sector hídrico al taller. \\
10 oct.\ 2023* & Taller participativo en el Centro Gaucho El Palenque: 88 personas acreditadas. El consultor excluye La Calderilla y la margen izquierda por los «alcances de la contratación»; los vecinos piden suspender la extracción, los desmontes y la venta de lotes. \\
25 oct.\ 2023* & Taller participativo de revisión del OTBN provincial, con carácter vinculante: 53 instituciones del Consejo Asesor y 91 participantes acreditados producen, en cuatro grupos, el único insumo del mapa final. \\
nov.\ 2023* & Peritaje judicial sobre el barrio privado Los Maitines en San Lorenzo (zona amarilla, Cat.\ II, Ley 26.331). \\
17 nov.\ 2023* & Ordenanza 607/23 de Vaqueros: modificaciones al Código de Planeamiento Urbano Ambiental y al Código de Edificación. \\
30 nov.\ 2023* & Informe final del equipo técnico del OTBN (SAyDS, INTA, INENCO-CONICET, UNSa, INAI, APN). \\
10 dic.\ 2023* & Decreto Nº 64/23: designa a la Dra.\ María Silvina Abilés Secretaria de Tierras y Bienes del Estado. \\
2024 & El diagnóstico social del Plan registra que el \textbf{bosque ribereño se eliminó durante la construcción del acueducto Campo Alegre}, con plan de forestación compensatoria del que quedan carteles. Y que los \textbf{tres puentes} del departamento son ``antiguos, obsoletos, sin mantenimiento y muy estrechos''. \\
2024 & \textbf{Febrero a julio}: la consultora releva el diagnóstico social con encuestas a vecinos, \textbf{bajo anonimato por temor a represalias}. YSATI Ingeniería presenta el Plan ---234 hojas y 12 planos---; el \textbf{14 de noviembre} la Dirección General de Infraestructura Hídrica no formula objeciones. \\
26 ene.\ 2024* & B.O.\ 21.640: la Provincia contrata por vía abreviada (art. 15 inc. e, Ley 8072) la consultoría externa que elaborará el Plan de Manejo Integral. Proveedor: YSATI Ingeniería S.R.L. Expte. 136-221704/23. La carátula nombra a Jardín Celestial S.R.L.\ como tercero y consigna el Expte.\ Judicial 238670/20 (probablemente el 23860/20, con un dígito de más). \\
18 feb.\ 2024* & Lluvias nocturnas anegan El Balcón, El Durazno y El Nogalar; el arroyo Guaranguay crece. \\
3 mar.\ 2024* & Temporal y desborde del río en el loteo Santiago Apóstol; el municipio asiste a cinco familias. \\
12 mar.\ 2024* & El Comité de Emergencia municipal entrega a los concejales un protocolo reformado para aprobar por ordenanza; los bomberos voluntarios informan que no tienen movilidad. \\
ago.\ 2024 & El municipio intima a propietarios con domicilio fiscal en otras localidades por morosidad en tasas e impuestos. \\
dic.\ 2024 & Comercialización activa de Solanas de La Caldera por Cabrera Bienes Raíces (material publicado). \\
dic.\ 2024* & El Senado aprueba la actualización del OTBN --- Ley Nº 8483 --- con tres votos en contra: Cachi, San Martín y \textbf{La Caldera}. \\
6 y 10 dic.\ 2024* & Se acompaña el \textbf{Plan de Manejo} al expediente de ejecución. \\
3 ene.\ 2025* & Se publica la Ley Nº 8483: revisa el OTBN e introduce el Área de Producción y Conservación, dentro de la cual las áreas verdes carecen de ubicación geográfica inicial. Su Anexo I asigna a la cuenca Mojotoro--Lavayén --- que contiene la microcuenca del río La Caldera --- 29.744 ha de APC y \textbf{0,0\% categorizable a verde}: por el criterio de la propia norma, la deforestación queda prohibida en esa cuenca, que la Tabla 17 del informe sombrea como excedida. Según las capas del informe, la cuenca cubre tres cuartos del departamento. \\
6 feb.\ 2025* & Decreto Nº 63/25: adjudica en venta ciento tres parcelas fiscales del barrio Juan Manuel de Rosas, \textbf{departamento Capital}, por adjudicación directa a actuales ocupantes (art. 2 inc. 1, Ley 2616). Expte.\ 1100388-113156/2024-0. \\
13 mar.\ 2025* & El proyecto de reserva \textbf{caduca} por aplicación de la Ley 1111, sin haber sido tratado. \\
20 mar.\ 2025* & La \textbf{Resolución 000056/2025} determina la línea de ribera del \textbf{arroyo Vaqueros} a la latitud del catastro 1122, y la del río Ancho en Cerrillos. \textbf{Los puntos de vinculación se remiten a la sede del organismo}: diez años después, la práctica de no publicar coordenadas sigue en uso. \\
abr.\ 2025 & Acto de entrega en Juan Manuel de Rosas. La presidenta del centro vecinal cuestiona la atribución del logro y sostiene que 26 de 30 adjudicatarios tienen domicilio en La Caldera. \textit{Imputación no verificada.} El acto de abril no coincide en fecha ni en cantidad con el Decreto 63/25. \\
9 may.\ 2025* & Res.\ 108/2025 de la Secretaría de Recursos Hídricos: rectifica las determinaciones 279/17 y 306/19 del río La Caldera, matrícula 4146. Los puntos se remiten a la sede. \\
11 may.\ 2025 & Elecciones provinciales. Daniel Roberto Moreno Ovalle\index{Moreno Ovalle, Daniel Roberto}, intendente de Vaqueros, es electo senador por La Caldera con 2.989 votos (27,94\%) contra 2.830 (26,45\%) de Héctor Miguel Calabró\index{Calabró, Héctor Miguel}, senador en ejercicio, según el escrutinio definitivo. Los porcentajes son sobre 10.698 votos válidos; se emitieron 11.171 sobre un padrón de \textbf{17.068 electores}, en un departamento con 12.299 habitantes y 8.976 mayores de dieciocho años censados en 2022. \\
30 jun.\ 2025* & Res.\ 142/2025: determina la línea de ribera del arroyo Quintín en el catastro 3756, casi doce años después de la comisión de la Res.\ 292/13. \\
3 jul.\ 2025* & Ante la denuncia de un diputado sobre «devolución de obras», el intendente y el ministro de Infraestructura explican que un convenio de defensas de 2023 (\$63 millones) se rescindió por desactualizado. \\
27 ago.\ 2025* & Resolución SAyDS Nº 00509: habilita el cambio de uso de suelo en inmuebles con APC, identificando el titular la Categoría III dentro del área. \\
oct.\ 2025* & Prórroga de la ley de suspensión de desalojos rurales hasta el 31/12/2026, a iniciativa del senador por La Caldera. \\
20 oct.\ 2025* & Disposición nacional Nº 2483/2025: acredita la actualización del OTBN salteño. \\
21 oct.\ 2025* & Res.\ ENARGAS 778/2025: pone en consulta pública el proyecto que reemplazará el régimen de expansión de redes de gas. \\
nov.\ 2025 & Según la prensa ---la lista de actuaciones no lo registra---, fallo de ejecución: participación vecinal en el manejo de la cuenca, línea de base socioambiental y reparación por manejo irregular de áridos. \\
nov.\ 2025 & Sesión de incorporación del Senado provincial. El mandato de Calabró cesará el 10 de diciembre; es el único representante departamental que votó contra la Ley 8483. \\
6 nov.\ 2025* & Res.\ ENARGAS 875/2025: amplía el plazo de la consulta. \\
27 nov.\ 2025* & \textbf{Naturgy NOA S.A.}, licenciataria de distribución sobre Salta, presenta observaciones en la consulta (Expte.\ EX-2025-112087756-APN-GD\#ENARGAS). \\
10 dic.\ 2025 & Moreno renuncia a la intendencia de Vaqueros y asume la banca. Ejerce el Ejecutivo municipal el concejal Daniel Viveros, presidente del Concejo (art.\ 64 de la Ley 8126; el art.\ 67 manda convocar a elecciones si falta más de un año de mandato). \\
28 dic.\ 2025* & Crecida violenta del arroyo Guaranguay. El amparista denuncia el incumplimiento de la cautelar. Al día siguiente, Recursos Hídricos presupuesta una obra de emergencia de \$27,8 millones. \\
ene.\ 2026 & Crecidas. El río socava el camino pasando El Durazno; decenas de familias aisladas. \\
3--4 ene.\ 2026* & Inundaciones en Nogalar 4, El Durazno y Potrero de La Caldera. El 5, el municipio pide \$10 millones a la Secretaría del Interior y maquinaria a Obras Públicas. \\
4 ene.\ 2026 & Cede el ingreso al Loteo Cerrado Solanas de La Caldera, con daños en media calzada de la RN 9, km 1618. \\
5 ene.\ 2026 & Obra provincial: badén sobre el arroyo Guaranguay, \$120 millones, expte.\ 366-16303/25, contratista Municipalidad de La Caldera. Se consigna que El Nogalar 1--4 son loteos municipales. \\
6 ene.\ 2026* & La jueza de feria intima a la Municipalidad y a la Provincia a informar en un día y ordena cumplir el Protocolo para Crecidas e Inundaciones, bajo apercibimiento de desobediencia judicial. \\
16 ene.\ 2026 & \textit{(Dato de prensa referido a la Capital, no al departamento.)} La Provincia deja sin efecto adjudicaciones de terrenos fiscales en la Capital por falta de ocupación efectiva y dispone reasignación por sorteo. \\
18 ene.\ 2026* & Nuevo evento: los sedimentos anegan los badenes de los arroyos de La Caldera. \\
19 ene.\ 2026 & La crecida del río Vaqueros arrasa una parte de la obra del nuevo puente que unirá Vaqueros con Salta. Más de treinta personas rescatadas por crecidas en La Caldera, Vaqueros y Campo Quijano. \\
21 ene.\ 2026* & Sesión Extraordinaria Nº 01 del Concejo: prórroga de la Ord.\ 733; clausura preventiva del ingreso a Solanas; clausura de dos campings; partidas para obras de defensa. Promulgación de las Ordenanzas 830, 831 y 832. \\
21 ene.\ 2026 & Reunión en Casa de Gobierno con el jefe de Gabinete y el ministro de Gobierno. El intendente propone un plan integral de manejo de cuenca; el jefe de Gabinete pide más control de áridos y que los municipios incrementen el control a los loteos. \\
21 ene.\ 2026* & Entre las promulgadas ese día, la \textbf{Ordenanza Nº 830} ---numerada /25 por su año de sanción--- prohíbe todo cobro de acceso al perilago de Campo Alegre. \\
21--27 ene.\ 2026* & Aguas del Norte y un equipo privado encauzan el río aguas arriba del puente para reparar el acueducto Norte, que había colapsado en Zabaleta; el 27 empiezan las defensas de Santiago Apóstol. \\
22 ene.\ 2026* & Inmuebles inscribe el embargo sobre la matrícula 1.846 de la Municipalidad ---donada por la Ley 7699--- en la ejecución de honorarios del amparo. \\
26 ene.\ 2026* & Ordenanza 648/26 de Vaqueros: modificaciones al Código de Edificación y de Planeamiento Urbano «año 2025», promulgada el mismo día por el intendente interino Daniel Viveros (Decreto 486/26). De ella resulta el texto del Código que analiza el capítulo \ref{cap:vaqueros} (art.\ 14: procesos de gentrificación). \\
27 ene.\ 2026 & El centro vecinal de La Calderilla denuncia desmonte de monte ribereño sin autorización y una audiencia pública sobre un loteo con sanción ambiental vigente. \\
2 feb.\ 2026* & Parte municipal de La Caldera: maquinaria pesada aportada por organismos provinciales ejecuta bateas, encauzamiento y muros de contención en el río La Caldera. Diez días antes de la audiencia. \\
12 feb.\ 2026* & Audiencia ante la jueza Mosmann con intervención del fiscal civil Agustín Vidal (Expte.\ 780280): se aprueba el Plan de Manejo Integral; se intima al Municipio a cumplir el Protocolo de Crecidas e Inundaciones a través del Comité Operativo de Emergencia \textbf{bajo apercibimiento de remitir los antecedentes a la justicia penal}; se fijan plazos a la Provincia para la planificación de obras; se ordena poner el Plan y la sentencia a disposición de la comunidad en el domicilio del colectivo actor; se reconoce la labor del COPAIPA. \\
18 feb.\ 2026* & Parte municipal: Tasa Neta de Ocupación Turística del 80,9\% en Carnaval, con cuatro establecimientos al 100\%. \\
mar.\ 2026 & La Corte de Justicia de Salta rechaza la impugnación y confirma a Salvatierra. Tercer titular del Ejecutivo municipal en tres meses. \\
1 mar.\ 2026 & El Concejo de Vaqueros elige presidenta a Andrea Salvatierra, que asume la intendencia interina. Viveros impugna. \\
20 mar.\ 2026* & Obras Públicas informa que las obras del Plan están a nivel de prefactibilidad, actualiza su costo a \$32.900 millones y admite que no hay presupuesto para los proyectos ejecutivos; reconoce asentamientos dentro del cauce. El 27, Fiscalía lo presenta al juzgado. \\
13 abr.\ 2026* & Resolución SAyDS Nº 00235: reglamento de requisitos para proyectos de intervención sobre bosques nativos. \\
15 abr.\ 2026* & Obras Públicas manda proyectar por administración las defensas del río y de los arroyos, «no contemplado en el Presupuesto 2026». \\
16 abr.\ 2026* & Último movimiento registrado del incidente. La causa continúa activa siete años y cuatro meses después de los hechos. \\
27 abr.\ 2026* & \textbf{Res.\ ENARGAS 435/2026}: aprueba el nuevo régimen de expansión de redes y deja sin efecto la Res.\ I-910/09. Reduce el horizonte de valuación de 35 a 10 años. \\
6 may.\ 2026* & Obras Públicas pide a Fiscalía una prórroga de sesenta días para la planificación de obras. \\
jun.\ 2026 & Cobertura sobre litigios de consumidores contra el grupo desarrollador y sobre la doble actuación de su abogado. \\
13 jul.\ 2026* & B.O.\ 22.224: Res. Delegada 567D/26. La Secretaría de Tierras y Bienes del Estado detecta errores materiales en los anexos de diecisiete decretos y una resolución delegada de adjudicación de tierra fiscal dictados entre 1991 y 2025, y los subsana en un solo acto. La mayoría son de grafía de nombres; dos casos (Leopoldo Lugones y Los Profesionales) cambian parcela y matrícula, dos corrigen el documento de identidad del adjudicatario y uno (Barrio Martel, Tartagal) convierte el precio en «sin cargo». \\
3 sep.\ 2026* & La \textbf{Resolución 479} (B.O.\ Nº 22.266, 10/09/2026) da de baja el contrato de defensas sobre el río La Caldera, de \$216.819.417, \textbf{porque cambiaron los niveles y el cauce}. No se precisa qué proyecto la reemplaza. Es lo mismo que constató el decreto de 1940. \\
4 sep.\ 2026* & Última actuación registrada en la ejecución del amparo; \textbf{la causa sigue abierta}. El Plan de Manejo se había aprobado el 12 de febrero, catorce meses después de presentado. \\
\bottomrule
\end{longtable}

\section*{Sesgos declarados}

\begin{enumerate}
\item \textbf{Sesgo de conflicto.} La cobertura se concentra en aluviones, audiencias y sentencias, y desaparece durante los períodos de tramitación administrativa, que son donde se decide.
\item \textbf{Sesgo de fuente.} Predominan medios con línea editorial crítica del desarrollo inmobiliario. La versión de los desarrolladores aparece casi exclusivamente en declaraciones de sus abogados. Falta relevar la prensa sectorial inmobiliaria.
\item \textbf{Sesgo de municipio.} La Caldera está mejor cubierta que Vaqueros, porque el litigio principal es caldereño.
\item \textbf{Densidad desigual.} El barrido de prensa fue pobre entre 2022 y 2024; las entradas de ese período provienen sobre todo de boletines y del expediente judicial. Corresponde barrido dirigido sobre actas municipales.
\end{enumerate}
